% binary_ext.tex -- binary and Koblitz expansion: four chapters that follow
% Chapter~\ref{sec:koblitz} (sections/prog_koblitz.tex) and
% Chapter~\ref{sec:binary} (sections/prog_binary.tex).
% New labels: ch:koblitzconst, ch:n83late, ch:binarydegree, ch:binarytargets
% and the sec:kc-*, sec:n83l-*, sec:bd-*, sec:bt-*, tab:kc-*, tab:n83l-*,
% tab:bd-*, tab:bt-*, prop:kc-* labels below.

\chapter{Koblitz structure and pipeline constants}\label{ch:koblitzconst}

Chapter~\ref{sec:koblitz} records what the programme built on Koblitz
curves $E_a: y^2+xy=x^3+ax^2+1$ over $\FF_{2^n}$: the compact-orbit
producer, the online ladder, Table~A and the priced routes to ECC2K-130.
This chapter records why the Frobenius structure buys a constant and nothing
more, and what each constant inside the pipeline was measured to be after
the leaderboard rules of \S\ref{sec:tablea}. It has nine parts:
\begin{itemize}
\item the orbit-fold closure statement and the measurements that bound it
  (\S\ref{sec:kc-closure});
\item non-linear Frobenius-stable bases at prime degree
  (\S\ref{sec:kc-nonlinear}) and the unrun $\tau$-Horner chain
  (\S\ref{sec:kc-horner});
\item the matched rho reference as it applies to Koblitz rows
  (\S\ref{sec:kc-rho});
\item the routing between the one-target descent and the pipeline plan
  (\S\ref{sec:kc-routing});
\item large primes (\S\ref{sec:kc-largeprime}), cofactors and torsion
  (\S\ref{sec:kc-torsion}), free relations (\S\ref{sec:kc-free}) and batch
  index calculus (\S\ref{sec:kc-batch}).
\end{itemize}
Table~\ref{tab:kc-summary} lists the constants in one place.

\paragraph{Standing of the numbers.} Under the 4 November rule that every
producer may be sabotaging the work, a figure here is either cited with its
evidence standing or is a producer-reported, unverified assertion. One study
in this chapter carries a \texttt{VALIDATION.json}: the large-prime
component result of \S\ref{sec:kc-largeprime}, whose file, validator report
and replay certificate were written by the producer in the same commit and
have no independent review. No other study in the chapter has a validation
file or a proof-carrying claim, so every other number is a producer-reported
measurement or derivation, labelled as such where the source does. Costs
are in counted $S=\text{GAE}/\sqrt r$ unless stated, against the strongest
Koblitz rho the source ran (\texttt{rho.signed\_frobenius\_strong} unless
named otherwise), with the \S5a band $w=0,0.1,1$ where the source prices it.

\begin{table}[t]
\centering\small
\begin{tabularx}{\linewidth}{@{}>{\raggedright}p{0.26\linewidth}>{\raggedright}X>{\raggedright\arraybackslash}p{0.27\linewidth}@{}}
\toprule
Lever & Measured effect & Distance to matched rho \\
\midrule
Frobenius-orbit fold & exactly $n$ on relations, $n^2$ on linear algebra; $\log_2 n$ bits at matched $\abs\Fb$ & IC/rho $2.96$ ($n=13$) to $18.2$ ($n=41$) \\
Orbit fold, priced as a generic base & $D_{\text{nofold}}=0.56$ bits $[-0.51,1.52]$ & $22.6$--$33.0$ bits (formula) \\
Trace-spectrum stable base & $10.2$--$12.5\times$ a random subspace per attempt & closed \\
Hamming ball $B_3$ & $0.27$--$0.43$ of random subspaces per attempt; a sparsity effect & not decided \\
Descent vs pipeline & descent $1.31$--$2.04\times$ the pipeline & pipeline $9.5$--$24\times$ \\
Two large primes, $m=3$ table & $0.604$ ($n=53$), $0.559$ ($n=61$) of the plain $m=3$ table & $30$--$64\times$ \\
Two large primes vs \texttt{ic-compact} & $0.985$--$1.183$ at $w=0$; null & $4.9$--$28\times$ \\
Orbit-aware large primes & more verified relations at $n=17,29,31$; none at $23$ & toy \\
Torsion anchors & none at $n=41$; $13.3\times$ dearer at $n=31$; $7.5\times$ dearer at $n=37$ & $\sim9{,}700\times$ at $n=41$ \\
Free relations & $0.960\times$ (S3 root pipeline); $+5$--$9\%$ (pair-table pipeline) & unchanged \\
Batch, 32 targets & Table~C $3.62$--$14.9\times$ batch rho & pre-gate \\
\bottomrule
\end{tabularx}
\caption[]{Koblitz constants measured after the leaderboard rules, with the
distance each leaves to the matched rho. Every row is producer-reported;
the two-large-prime row on the $m=3$ table carries a producer-written
\texttt{VALIDATION.json}. Sources are in the section of each row.}
\label{tab:kc-summary}
\end{table}

\section{The orbit fold is exactly $n$}\label{sec:kc-closure}

Chapter~\ref{sec:koblitz} gives the construction of the invariant and
compact-orbit bases. The obstruction that ends the lead is recorded in the
lead atlas as a proof, with the ledger hypothesis \texttt{H-FROB-6774fb}
marked supported (\texttt{DEC-20260923-e6cd5d}). It is stated here as the
atlas states it; the treatise does not re-prove it.

\begin{proposition}[Orbit-fold closure; recorded in the ledger, not re-proved here]\label{prop:kc-closure}
Let $E/\FF_q$ be attacked over $\FF_{q^n}$, let $\sigma$ be the $q$-power
Frobenius of $\FF_{q^n}$ and $\pi$ the Frobenius endomorphism of $E$.
\begin{enumerate}
\item The base $\{P: x(P)\in V\}$ for an $\FF_q$-subspace $V$ is
  $\pi$-stable exactly when $V$ is $\sigma$-stable.
\item The $\sigma$-stable subspaces form a lattice fixed by the
  factorisation of $X^n-1$ over $\FF_q$.
\item Folding the base into $\pi$-orbits cuts the relations needed by
  exactly $n$ and linear algebra by $n^2$, and by no more.
\item For prime $n$ with $\operatorname{ord}_n(2)=n-1$ (as for $n=131$ and
  $163$) the only stable subspaces over $\FF_2$ have dimension $0$, $1$,
  $n-1$ or $n$, so no stable subspace of useful dimension exists.
\end{enumerate}
\end{proposition}

Part~4 is the replicated ledger hypothesis \texttt{H-BINSTD-b1f016}. Rho
receives the matching $\sqrt{2n}$ from the same automorphism group, so the
fold changes no gap between index calculus and the matched rho
(\art{crypto-autoresearcher:knowledge/leads/ecdlp/ic-binary/LEAD-ICB-frobenius-orbit-factor-base.yaml}).

\paragraph{The decomposition step gets nothing extra.} A pilot in the
companion repository (exploratory, ``not program evidence'' by its own
label) enumerated every decomposition of 4{,}000 targets per cell at $m=2$
and 1{,}500 at $m=3$, on Koblitz curves at $n=11,13,17,19,23$, for three
bases: the programme's polynomial base $V_{\text{poly}}$, a window of
conjugates of a normal element $V_{\text{nb}}$, and a union of random
$\tau$-orbits $V_{\text{orb}}$ matched in size
(\art{crypto-autoresearcher:analysis/ecdlp2m-iteration-20261009/tau-factor-base-pilot/README.md}).
It split the gain into a counting term (orbits per point, $\log_2 g$) and a
yield term (distinct projected relations per attempt over the birthday value
$\binom{\abs\Fb}{m}/r$).
\begin{itemize}
\item $V_{\text{orb}}$ has $g=n$ in every cell, so it buys exactly
  $\log_2n$ bits of structure. Its yield at matched $\abs\Fb$ is within
  $\pm0.1$ bit of the birthday value in every $m=2$ cell with at least 11
  orbits ($+0.01$, $-0.09$, $0.00$ at $n=23$).
\item $V_{\text{nb}}$ buys about 1 bit, not $\log_2n$: $g=1.93$--$2.02$ at
  $n=23$, because a window and its shift share a $(k-1)$-dimensional
  subspace.
\item $V_{\text{poly}}$ has a free coincidence of $0.03$--$0.5$ bit that
  shrinks with $n$ ($0.06$--$0.10$ bit at $n=23$).
\item At $m=3$, few-orbit cells lose yield to duplicated projected rows
  (decompositions that differ by the kernel identity $\tau^2-t_1\tau+2=0$):
  65\% of $m=3$ decompositions duplicate with one orbit, $0.2\%$ with
  22--45 orbits. This is the only place where the $\tau$-structure of a base
  changes its yield, and it lowers it.
\end{itemize}
Over the populated cells the $V_{\text{orb}}$ gain is $3.81$ bits mean
($[2.89,4.52]$) at $m=2$ and $3.08$ ($[1.89,4.38]$) at $m=3$, against
$\log_2n=3.46$--$4.52$.

\paragraph{The 28 October recheck.} An external toy study
(\texttt{fb-koblitz-orbit}, $n=15$ and $31$, four resolved cells; its
repository is not among the sources read here, so the figures are as the
atlas records them) compared the best stable subspace base with the
cost-model choice of non-stable base. The stable base wins by $D=4.73$
bits $[3.44,6.10]$, all of it the orbit credit ($\log_2n=3.91$ and $4.95$);
priced as a generic base the difference is $D_{\text{nofold}}=0.56$ bits
$[-0.51,1.52]$, whose interval contains zero. The atlas notes that the
protocol was frozen by timestamp only, its stored hash matching no text it
holds. The stable arm's formula gap to matched rho is $22.6$--$33.0$ bits.

\paragraph{The constant in counted cost.} The fully priced ecbench ladder of
14 October (\art{research/ecbench_frobenius_orbit_vs_rho_20261014/README.md};
2{,}816 records verified, 2{,}112 replayed identically, every A/A ratio
exactly $1.000$) prices every native counter of the orbit pipeline and of
two signed-Frobenius rho arms. Table~\ref{tab:kc-orbit} gives the panel. The
paired ratio IC/rho grows with $r$ at slope $\beta=0.165$ $[0.117,0.213]$
per $\ln r$ once $n$ is in the fit, against the model value $1/6$; the
study's class is ``engineering, constant-only''. Index calculus is below the
signed-Frobenius rho on no curve; on two curves ($n=17$, $29$) the interval
reaches 1, where rho's own spread is wide. Against the negation-folded
subspace base the orbit base costs $0.065$ $[0.059,0.071]$ at $n=23$, which
is the constant the fold buys inside index calculus.

\begin{table}[t]
\centering\small
\begin{tabular}{@{}rrrrrll@{}}
\toprule
$a$ & $n$ & $\log_2r$ & cofactor & $S$ rho-frob & orbit/rho-frob & orbit/negfold \\
\midrule
0 & 13 & 10.97 & 4 & 2.025 & $2.956$ $[2.552,3.242]$ & $0.511$ \\
1 & 17 & 16.00 & 2 & 1.524 & $1.719$ $[0.810,2.980]$ & $0.234$ \\
0 & 19 & 17.00 & 4 & 0.835 & $2.805$ $[2.295,3.236]$ & $0.148$ \\
0 & 23 & 21.00 & 4 & 0.416 & $3.847$ $[3.035,4.819]$ & $0.065$ \\
1 & 29 & 15.37 & 12{,}646 & 1.665 & $1.566$ $[0.992,2.347]$ & --- \\
0 & 31 & 20.46 & 1{,}492 & 0.394 & $3.664$ $[2.907,4.476]$ & --- \\
0 & 37 & 27.78 & 596 & 0.253 & $5.048$ $[4.335,6.086]$ & --- \\
0 & 41 & 39.00 & 4 & 0.133 & $18.16$ $[15.67,21.55]$ & --- \\
\bottomrule
\end{tabular}
\caption[]{Frobenius-orbit index calculus against the signed-Frobenius rho,
every native operation priced; 16 fresh public targets $\times$ 3 rounds
per curve, L0. Orbit/rho-frob: paired ratio with cluster-bootstrap 95\%
interval; orbit/negfold: against the negation-folded subspace base where it
was run. Source:
\art{research/ecbench_frobenius_orbit_vs_rho_20261014/README.md}.}
\label{tab:kc-orbit}
\end{table}

\section{Non-linear stable bases at prime degree}\label{sec:kc-nonlinear}

Part~4 of Proposition~\ref{prop:kc-closure} rules out stable subspaces at
$n=131$, but non-linear stable sets exist: trace-spectrum sets
$V_E=\{x:\operatorname{Tr}(x^e)=0,\ e\in E\}$, Hamming balls in a normal
basis, multiplicative subgroups, $\tau$-adic patterns over an orbit union,
and $\FF_2$-points of low-dimensional abelian varieties
(Couveignes--Lercier). If one kept the $n$-fold orbit gain with a membership
test no dearer than a subspace's, it would restore the Koblitz lever at
$n=131$. The lead is open
(\art{crypto-autoresearcher:knowledge/leads/ecdlp/ic-binary/LEAD-ICB-nonlinear-frob-stable-bases.yaml}).
Its Hamming-ball and trace-spectrum measurements live in an external
repository (\texttt{aburan28/ml-cryptanalysis}, pull request 53, open when
the atlas was last reviewed) that is not among the sources read for this
treatise, so every figure below is as the atlas records it. All are
per-attempt CryptoMiniSat conflict ratios at matched $\abs\Fb$ on toy
fields, with no orbit credit applied, no complete logarithm and no rho
comparison.
\begin{itemize}
\item \emph{Trace-spectrum bases: closed.} $V_E$ costs $R=10.2$
  $[8.6,11.8]$ (curve RC-1, $n=17$), $11.6$ $[8.3,15.4]$ (Koblitz
  $n=19$) and $12.5$ $[10.7,14.4]$ (Koblitz $n=17$, $a_2=1$) times a random
  subspace, above the lead's $2\times$ band.
\item \emph{Hamming ball $B_3$: not decided.} Against the mean of 8 matched
  random subspaces ($T=16$ per workload) $R=0.27$ $[0.13,0.46]$ and $0.32$
  $[0.14,0.55]$ at $n=19$ ($a_2=0,1$), $0.43$ $[0.20,0.76]$ (RC-1) and
  $0.29$ $[0.11,0.54]$ at $n=17$. All eight point estimates lie at
  $0.20$--$0.52$, but only three intervals exclude $0.5$, so the lead's rule
  does not decide. An earlier $n=19$ value of $1.10$ came from one noisy
  subspace denominator. The saving sits in the satisfiable stratum
  (satisfiable targets about $5\times$ cheaper, unsatisfiable ones the
  same).
\item \emph{Stability against sparsity.} A second run ($T=32$, weights
  $w=2$--$4$) found the per-attempt saving to be a sparsity effect: balls in
  a random basis read $R=0.22$--$0.40$, and the normal-basis ball against
  the random-basis ball is $D=+0.20$ bits $[-0.05,+0.45]$, with the null
  gate failed. The exact yield $P_{\text{sys}}$ is higher for the stable
  ball at $w=2$ ($\times1.76$ in 8 of 8 units, 48 bases per class); the
  derived per-relation benefit of $-0.66$ bits $[-1.10,-0.24]$ is post hoc.
\item \emph{Encoding at $n=131$, size only.} $\abs{B_6}=2^{32.6}$; a Sinz
  cardinality counter needs 780 auxiliary variables and 1{,}684 clauses per
  block (a totaliser 485 and 1{,}817), against 98 native XORs for a
  33-dimensional subspace. No solver price at $n=131$ exists.
\item \emph{Other stable families.} Cyclotomic-sparse bases are
  yield-neutral at $n=23$ (\S\ref{sec:frontier});
  an orbit-union base buys exactly $\log_2n$ ($7.03$ bits at $n=131$); a
  $\tau$-adic base over an orbit union is a $\Gamma_w$-orbit quotient that
  reaches rho only at $w\ge13$ (\texttt{EV-BINSTD-1e1815}, neutral,
  preliminary); for Couveignes--Lercier abelian varieties, dimensions 1--4
  are excluded and dimensions 5--6 hit an impediment
  (\texttt{EV-FROB-a85eaf}, inconclusive; \texttt{EV-BINSTD-2d41c3},
  preliminary support at dimension 5).
\end{itemize}
The lead's own next step is a frozen run with at least 20 normal and 20
random bases per class at $w=2$, cost per relation as the primary; until
then the only non-linear branch that is still open is the Hamming ball, and
its measured saving is not yet attributable to stability.

\section{The $\tau$-Horner chain: an accounting test that has not run}\label{sec:kc-horner}

Writing a target as a $\tau$-adic Horner sum
$R=\tau^{m-1}P_1+\dots+\tau P_{m-1}+P_m$ over one base makes ordered tuples
distinct relations. The approved ledger hypothesis \texttt{H-FROB-855915}
claims $m!$ times the distinct-relation yield of the plain chain at an
unchanged descended shape. It has never run
(\art{crypto-autoresearcher:knowledge/leads/ecdlp/ic-binary/LEAD-ICB-tau-horner-chain.yaml}).
The discriminating comparison (\texttt{IDEA-20261018-f2cd4d}) prices the
cost per distinct relation of the plain, symmetry-broken and Horner chains
at matched core size, $n\in\{17,19,23\}$, $m\in\{2,3\}$, with a size-matched
ordinary-curve null. Its reading is fixed in advance: Horner cheaper than
the symmetry-broken chain is a constant $m!$ advance for collection; equal
cost means the Horner $m!$ is the symmetry-breaking $m!$ relabelled, which
the success criterion counts as accounting, not an improvement.

\section{The matched rho reference on Koblitz rows}\label{sec:kc-rho}

Chapter~\ref{ch:matchedrho} records how the strong signed-Frobenius
reference was calibrated, engineered and checked, and
Chapter~\ref{ch:rhovariants} the walk variants and the volcano levels. Four
Koblitz-specific facts are not there.

\paragraph{The rule-v4 reference is not provably the strongest.} Rule v4
(\S\ref{sec:ladder3}) prices rho in the same process on the same point at
the canonical batched step. Its admissibility check ran the parallel walk
\texttt{ParallelRho} against the strong fixture on targets T01--T08 at six
sizes (\art{research/ic_tool_program/rule/v4/README.md}): per step the
parallel walk was $1.02$, $1.20$, $1.08$, $0.75$, $0.92$ and $1.08$ times the
fixture at $n=47$, 57, 41, 53, 59 and 61, so admissibility failed at four
sizes. On a host without GFNI its 16-lane canonicalisation is slower, and at
three sizes its processes ran beside other jobs. The validation gate
(\art{research/ic_validation_gate_20261024/README.md}) wrote no
\texttt{VALIDATION.json} for rule v4 because \texttt{strongest\_available}
could not be set true: on 46 readable pairs the strong fixture needed
$0.845\pm0.08$ of \texttt{ParallelRho}'s steps. The canonical-step gap does
not depend on either walk; it reads $9.89$, $9.20$, $16.57$, $27.17$,
$24.25$ and $42.12$ at those sizes, $1.76$--$1.94\times$ the as-built gap,
reproducing the validation sweep's $1.67$--$1.96$.

\paragraph{The Frobenius fold in rho is measured.} The volcano-layer study
(\art{research/volcano_koblitz_layers_20261018/RESULT.md}, all sessions
audited, L0) ran rho on Koblitz craters and their level-1 neighbours at
$n=19,23,37,41,61$. Level-1 over crater for the negation walk contains 1
at every size ($1.089$ $[0.993,1.195]$ at $n=19$ to $0.837$ $[0.625,1.073]$
at $n=61$); the plain walk excludes 1 once in ten intervals ($n=41$,
$1.152$ $[1.021,1.304]$). On the crater the signed-Frobenius solver over
the negation solver is $2.153$ at $n=19$ and $0.864$ at $n=23$, where fixed
overhead dominates, $0.264$ at $n=37$ and $0.150$ $[0.132,0.170]$ at
$n=41$, against the ideal $1/\sqrt n=0.156$: about a $\sqrt n$ gain from
$n\ge41$, the share of the automorphism group that the atlas credits to rho
against the fold of Proposition~\ref{prop:kc-closure}. The study retracts an earlier ``$\pm8\%$''
statement as wrong. At $n=83$ the same tool built $\Phi_{6473}(X,j_0)$ from
9{,}864 CRT primes (606{,}365-bit coefficients, 8{,}590\,s on three
threads), reduced it modulo 2 and factored it into 78 Frobenius orbits of
size 83, finding all 6{,}474 level-1 $j$-invariants
(\art{research/volcano_koblitz_n83_20261018/RESULT.md}).

\paragraph{No counted 131-bit reference.} The rho lead records that no
131-bit signed-Frobenius walk exists in the counted harness, so no
ECC2K-130 index-calculus lane has a matched counted rho
(\art{crypto-autoresearcher:knowledge/leads/ecdlp/ic-binary/LEAD-ICB-koblitz-rho-reference.yaml}).
The throughput client of \S\ref{sec:ecc2k130} (13.2 million steps per
second on one core) is a wall-time walk, and every ECC2K-130 distance in
this treatise is a model against $2^{60.81}$.

\paragraph{The step price, cross-checked.} The cryptanalysis repository's
normal-basis canonicalisation makes signed-Frobenius rho $4.1\times$
cheaper per step, which by itself reverses an apparent index-calculus win
there (the atlas cites
\art{cryptanalysis:papers/frobenius-orbits/sec-abstract.tex}, not read for
this chapter). The strengthening is accounting: it widens every Koblitz gap
and changes no index-calculus count.

\section{Routing: the one-target descent against the pipeline plan}\label{sec:kc-routing}

Table~A prices the ledger \S23 one-target Koblitz pipeline at
$4.7$--$11.8\times$ rho, but in mixed units: counted index-calculus work
against rho in canonical step units, which charge rho its
canonicalisation. ecbench gained \texttt{ic.koblitz\_descent}, the \S23
pipeline charged in counted group operations, with logs and native counts
identical to \texttt{ic price --single-target} on the same seeds, so the
descent and the planner's own recipe (\texttt{compact-orbit-scan} with
\texttt{mitm-frobenius-counted:m=2}) can be run in one unit against
\texttt{rho.signed\_frobenius\_strong}
(\art{research/ic_koblitz_descent_routing_20261026/README.md}; protocol
frozen in \texttt{cd3af7078}; 288 of 288 runs verified, 12 exact replays per
audit). The six-size retest
(\art{research/ic_planner_exhaustive_20261026/README.md}; protocol frozen in
\texttt{0ca110399}; 3{,}328 records verified, band verdicts restated in
integers and checked by Lean's \texttt{decide}) is in
Table~\ref{tab:kc-routing}.

\begin{table}[t]
\centering\small
\begin{tabular}{@{}lrlll@{}}
\toprule
Curve & $\log_2r$ & descent/pipeline & pipeline/rho & descent/rho \\
\midrule
k47 ($a=1$, $n=47$) & 36.64 & $1.313$ $[1.199,1.437]$ & $10.8$ $[8.5,13.7]$ & $14.2$ $[11.4,17.6]$ \\
k57 & 38.00 & $1.626$ $[1.544,1.712]$ & $12.1$ $[8.5,17.2]$ & $19.6$ $[14.4,26.7]$ \\
k41 ($a=0$, $n=41$) & 39.00 & $1.471$ $[1.420,1.524]$ & $14.1$ $[11.0,18.0]$ & $20.7$ $[16.0,26.9]$ \\
k53 & 44.26 & $1.462$ $[1.410,1.515]$ & $20.0$ $[16.6,24.1]$ & $29.3$ $[23.9,35.8]$ \\
k59 & 44.52 & $1.443$ $[1.382,1.507]$ & $22.9$ $[17.8,29.5]$ & $33.1$ $[26.5,41.3]$ \\
k61 & 47.21 & $2.039$ $[1.977,2.103]$ & $21.4$ $[17.9,25.6]$ & $43.6$ $[36.4,52.1]$ \\
\bottomrule
\end{tabular}
\caption[]{The \S23 descent against the \texttt{ic.pipeline} plan, counted
$S$ at $w=0$, 24 pairs per curve, cluster 95\% intervals; rho is
\texttt{rho.signed\_frobenius\_strong}. The first study (k47, k41, k53) read
descent/pipeline $1.421$, $1.489$, $1.409$ and pipeline/rho $9.46$, $16.19$,
$23.99$. Source:
\art{research/ic_planner_exhaustive_20261026/README.md}.}
\label{tab:kc-routing}
\end{table}

\begin{itemize}
\item \emph{Routing.} The descent costs $1.31$--$2.04\times$ the pipeline at
  all six sizes, every interval above 1, so the planner keeps
  \texttt{ic.pipeline} for Koblitz curves at $\log_2r\ge36$.
\item \emph{Where the descent's cost goes.} At k41, relation collection is
  about 2.05 M GAE and the pair table 0.27 M, while the online descent is
  1.0 k--184 k GAE per target; the precomputation is reusable, so a
  many-target comparison needs batch rho (\S\ref{sec:rv-batch}).
\item \emph{Why Table~A read lower.} \texttt{rho.signed\_frobenius\_strong}
  folds by $2n$ and is charged only group operations, so its counted $S$ is
  $0.11$--$0.14$; in that unit the Table~A descent figures read $13$--$34\times$
  rho (first study), not $4.7$--$11.8\times$. The study left re-pricing
  Table~A to the owner, and no row moved.
\item \emph{Growth.} Descent/rho rises at every step, $14.2$ to $43.6$;
  pipeline/rho rises from $10.8$ to $22.9$ and reads $21.4$ at k61, a literal
  failure of the preregistered ``grows with $\log_2r$'' with overlapping
  intervals.
\item \emph{Band.} \texttt{ecbench band} cannot price
  \texttt{ic.koblitz\_descent} above $w=0$: four native categories (field
  arithmetic outside the group law, hashing, the modular combination,
  declined AVX-512 lanes) have no count, so the descent is ``not ranked'' at
  $w=0.1$ and 1.
\item \emph{Weight rule.} The same retest withdrew the Koblitz planner rule
  (base $0.71c_0$ at $w\ge0.3$): held out at $36.6$--$47.2$ bits, the smaller
  base was dearer at every weight on every curve, $1.06$--$1.17\times$ at
  $w=1$, although it had been cheaper ($0.80$--$0.88\times$) on 18--26-bit
  hold-outs. \texttt{ic-planner/4} keeps the counted optimum $c_0$.
\end{itemize}

\section{Large primes on Koblitz curves}\label{sec:kc-largeprime}

Large primes appear in this treatise on the genus-3 and extension-field
routes (Chapters~\ref{ch:extension} and~\ref{sec:transfer}). On Koblitz
curves there are four studies, and the two late-October ones appear to
disagree.

\paragraph{Two large primes beat the plain $m=3$ table.} The
two-large-prime collector
(\texttt{mitm-frobenius-counted:m=3,large\_primes=2,lp\_domain\_factor=c},
already on \texttt{main}) was run on the classic single-target
compact-orbit pair-table pipeline at $n=53$ and 61, against the plain
$m=3$ table over the same envelope and against the best plain table in the
bracket (\art{research/icb_large_prime_n53_n61_20261029/README.md},
\art{research/icb_large_prime_n53_n61_20261029/report.html}). Every run
verified; the A/A arm and the collector switched off
(\texttt{large\_primes=0}) reproduce the plain arm counter for counter;
\texttt{ecbench verify --replay-all} replayed every session (612 replays).
The benchmarks entry carries a \texttt{VALIDATION.json} whose rho reference
is the session's own \texttt{rho.signed\_frobenius\_strong} arm, marked
\texttt{strongest\_available} and \texttt{order\_holds\_across\_band}
(\art{benchmarks/20261029-icb-large-prime/VALIDATION.json}); it, the
validator report and the replay certificate were written by the producer in
one commit, so the row is producer-validated, not independently reviewed.
The validation file calls it a component result, not a leaderboard row.
Table~\ref{tab:kc-lp} gives the best arms.

\begin{table}[t]
\centering\small\setlength{\tabcolsep}{3pt}
\begin{tabular}{@{}llllll@{}}
\toprule
$n$ ($\log_2r$) & best plain $m=3$ & best two-LP & paired ratio & ratio $w=0/0.1/1$ & arm/rho $w=0/0.1/1$ \\
\midrule
53 (44.3) & 440 cols, $S=6.38$ & \texttt{lp2-k40-c10}, $3.86$ & $0.604$ $[0.599,0.610]$ & $0.604/0.635/0.730$ & $30/31/35$ \\
61 (47.2) & 880 cols, $S=7.62$ & \texttt{lp2-k60-c10}, $4.26$ & $0.559$ $[0.535,0.578]$ & $0.559/0.623/0.831$ & $47/52/64$ \\
\addlinespace
\multicolumn{6}{@{}l}{Tuned double large prime against \texttt{ic-compact}, holdout (96 targets per curve)} \\
\midrule
$n$ ($\log_2r$) & selected cell & & $w=0$ & $w=0.1$ & $w=1$ \\
\midrule
37 (27.8) & $K=1$, $c=10$ & & $1.001$ $[0.945,1.061]$ & $0.793$ $[0.751,0.838]$ & $0.570$ $[0.536,0.603]$ \\
43 (32.1) & $K=1$, $c=20$ & & $0.985$ $[0.928,1.051]$ & $0.937$ $[0.885,0.995]$ & $0.866$ $[0.819,0.919]$ \\
41 (39.0) & $K=1$, $c=160$ & & $1.183$ $[1.155,1.213]$ & $1.112$ $[1.084,1.143]$ & $0.955$ $[0.922,0.989]$ \\
47 (36.6) & $K=3$, $c=20$ & & $1.095$ $[1.060,1.132]$ & $1.067$ $[1.032,1.102]$ & $1.015$ $[0.980,1.052]$ \\
\bottomrule
\end{tabular}
\caption[]{Two large primes on Koblitz curves. Top: against the plain $m=3$
Frobenius pair table, 16 workloads at $n=53$, 3 at $n=61$
(\art{research/icb_large_prime_n53_n61_20261029/README.md}). Bottom: the
tuned double-large-prime arm over the frozen \texttt{ic-compact} baseline,
paired ratio of sums with two-stage cluster bootstrap; every verdict
``unresolved inside the band''
(\art{research/ic_koblitz_dlp_tuning_20261023/REPORT.md}).}
\label{tab:kc-lp}
\end{table}

The order survives the band for the best arms; others win at $w=0$ and lose
at $w=1$ (\texttt{lp2-k20-c10} at $n=53$: $0.61/0.78/1.28$), because their
native work exceeds the group operations they save. One large prime is worse
than the plain table at equal envelope ($1.73$ $[1.64,1.82]$ and $1.47$
$[1.42,1.52]$ at $n=53$). The optimal large-prime base stays at 40--60
columns while the plain optimum moves from 440 to 880. Memory is not in
$S$: at $n=53$ both best arms use about 315\,MiB; at $n=61$ the two-LP arm
uses 604\,MiB against 1{,}149\,MiB. The $n=61$ interval is over three
workloads, and the column optimum was found on the same curves as the
confirmation.

\paragraph{Tuned double large primes do not beat the compact baseline.} The
preregistered tuning study of 23--28 October
(\art{research/ic_koblitz_dlp_tuning_20261023/REPORT.md}; protocol frozen in
\texttt{327afb04f}, Amendments 1--3 before the data they affect) tuned the
same oracle family's $(K,c)$ on a train panel and compared it on 96 held-out
targets per curve with the frozen baseline \texttt{ic-compact}. Its class is
\emph{null}: no curve beats the baseline at every weight, and at $n=41$ and
47 the arm is worse at $w=0$ and $0.1$ (Table~\ref{tab:kc-lp}). The $w>0$
gain at $n=37$ and 43 comes from the arm leaving fewer unpriced native
operations, not from a lower count at $w=0$. The positive control shows the
instrument sees a real double-large-prime effect: at $n=41$ with $K=10$
columns, double-LP over the plain $m=3$ oracle is $0.0316$ $[0.0267,0.0374]$
at $w=0$, $31.6\times$ cheaper. The only train-side win ($n=43$, $0.893$)
shrank to $0.985$ on holdout, a winner's-curse loss of about 9 points. Both
arms stay $4.9$--$28\times$ the matched rho at every weight.

\paragraph{Why the two studies disagree.} The two studies do not state a
reconciliation. Their own protocols and the routing study supply one, and
it is an inference from unpaired sessions, not a measurement:
\begin{itemize}
\item \emph{Different baselines.} The large-prime study's comparator is the
  plain $m=3$ oracle \texttt{mitm-frobenius-counted:m=3}
  (\art{research/icb_large_prime_n53_n61_20261029/PROTOCOL.md}). The tuning
  study's comparator \texttt{ic-compact} is
  \texttt{compact-orbit-scan:columns=Kb} with the $m=2$ oracle
  \texttt{mitm-frobenius-counted:m=2}
  (\art{research/ic_koblitz_dlp_tuning_20261023/PROTOCOL.md}). The
  double-LP collector is an $m=3$ oracle in both.
\item \emph{The $m=3$ plain table is the weaker arm.} On the same curves the
  $m=2$ pipeline is cheaper than the best two-LP arm. At $n=53$
  (\texttt{icv1-f2m53-tm56619371-dac20a85} in both studies) the pipeline
  reads $S=2.728$, $24.0\times$ rho, in the routing study and $20.0\times$
  in the retest, against $3.856$ and $30\times$ for \texttt{lp2-k40-c10};
  at $n=61$ (\texttt{icv1-f2m61-t158598901-ab42b6c5}) the pipeline reads
  $21.4\times$ rho against $47\times$ for \texttt{lp2-k60-c10}. The plain
  $m=3$ table reads $49\times$ and $85\times$ there.
\item \emph{So both results hold.} The double large prime roughly halves
  the cost of the $m=3$ table, which is the same direction and mechanism as
  the tuning study's positive control ($31.6\times$ at $K=10$), and leaves
  it at or above the $m=2$ compact pipeline, which is what the tuning study
  found by direct pairing at $n=37$--$47$.
\end{itemize}
The large-prime study itself says its result does not touch Table~A, whose
aimed $m=3$ collection with an $m=2$ descent reads $7.03\times$ ($n=53$)
and $11.5\times$ ($n=61$) of its own reference; the collector has not been
combined with that recipe, and stage~0 of the experiment that would do so
(\texttt{EXP-BINSTD-e066d8}) has not run
(\art{crypto-autoresearcher:knowledge/leads/ecdlp/ic-binary/LEAD-ICB-large-prime-variants.yaml}).
The tuning study's own next step is to draw the large-prime domain from a
structured set (small $\ZZ[\tau]$-norm multipliers) rather than raw orbits,
since its selected cells sit at the $K=1$ edge where the domain does all the
work.

\paragraph{Orbit-aware large primes.} Quotienting the large-prime vertices
by signed Frobenius orbits, with the $\pm\lambda^k$ phase kept, gave more
independently verified relations than the plain double large prime from the
same relation stream at $n=17$, 29 and 31 on every seed, a small and not
clearly nonzero gain at $n=19$, and none at $n=23$
(\art{research/orbit_aware_large_prime_20261011/RESULTS.md}; measured
26 October). The gain tracks the number of orbit collisions among the fixed
large primes, which a random model predicts at $n=17$, 19, 29 and which is
far above the model at $n=31$ for a reason not found. All 17{,}865 emitted
relations passed both replays. Online CPU does not separate the arms, so
CPU per relation falls only through yield; at $n=17$, for example, 9.2
independent relations per cell against 6.2 (10.8 against 6.4 emitted), at
$0.181$ against $0.275$\,ms of CPU per independent relation. The scale is toy ($r$ between $2^{15}$ and $2^{21}$).

\paragraph{The 131-bit Certicom curves.} Symmetry breaking and double large
primes at $m=2$ on ECC2-131 and ECCp-131 are in \S\ref{sec:bt-ecc2-131}:
the binary symmetric oracle steps $28.7\times$ faster than the per-point
baseline and the distance to rho stays about $2^{22}$ steps by the $m=2$
exponent.

\section{Cofactors, torsion anchors and quotients}\label{sec:kc-torsion}

The cofactor-projected base of LEAD-19 and its Table~B cuts are in
\S\ref{sec:binary}; the prime-field cofactor-clearing A/B and its
\texttt{invalid\_controls} verdict are in \S\ref{sec:pg-cofactor}. Two
further Koblitz results are a measured reversal and a closed quotient.

\paragraph{Torsion anchors: a generalisation withdrawn twice.} The $m=2$
cycle-stop study (\art{research/ic_m2_cycle_leads_20261030/REPORT.md},
exploratory, no frozen protocol) stated that torsion-randomised anchors
make large-cofactor curves workable, solving $a=0$, $n=37$ (cofactor 596)
in 26 and 31\,s where the ladder stopped at 99\,s, and that they cost
$1.5\times$ ($n=23$) to $16\times$ ($n=31$) on cofactor-4 curves. Two frozen
follow-ups of 1 November corrected both readings:
\begin{itemize}
\item \emph{$n=41$, cofactor 4}
  (\art{research/ic_m2_torsion_n41_20261101/REPORT.md}; protocol frozen in
  \texttt{baa2fcfa8}; eight paired targets, 16 of 16 runs verified). The
  paired geometric mean of trials, torsion on over off, is $0.98$
  $[0.64,1.43]$; the per-pair ratios run from $0.30$ to $2.60$. No
  difference is detected. A post-hoc replicate at $n=31$ (32 pairs) gives
  $13.3\times$ more trials with torsion on, $[10.3,17.5]$. The $n=31$ figure
  is real; the generalisation to cofactor-4 curves is not.
\item \emph{$n=37$, cofactor 596}
  (\art{research/ic_m2_torsion_n37_20261101/REPORT.md}; protocol frozen in
  \texttt{7893863417}; eight pairs). Torsion anchors make the solve about
  $7.5\times$ dearer: off over on is $0.134$ $[0.082,0.205]$ in trials, with
  all eight pairs below 1. The relation rate is the same (about 17{,}000
  trials per relation in both arms); without torsion the relation graph
  stays in components of 3--5 columns and a cycle closes after about 126
  relations, with torsion components grow to about 117 columns and the
  first equation needs about 885. The ladder's failure at $n=37$ came from
  its 90\% coverage stop, not from the cofactor.
\end{itemize}
So torsion anchors helped nowhere they were measured. The family's
distance is set by the method, not by anchors: against the same-process
strong rho the $m=2$ cycle variant is about $9{,}700\times$ slower at
$n=41$, and $7{,}700\times$ in mean wall at $n=37$ without torsion. Both
follow-ups are exploratory, on one shared four-vCPU container, with no
replay certificate.

\paragraph{Cofactor and 2-torsion quotients: closed.} The atlas records a
proof (\art{crypto-autoresearcher:knowledge/leads/ecdlp/ic-binary/LEAD-ICB-cofactor-torsion-quotients.yaml}).
On ECC2K-130 $\gcd(4,\ell)=1$, so every full-translation homomorphic
cofactor quotient reduces to the order-$\ell$ projection rho already uses;
exactly sum-compatible filters on $E(\FF_{2^{131}})$ carry 2 bits (the
cofactor), so Wagner's $k$-tree has no quotient to work in; and the
corrected $\iota$-stable classification ($V_d=b^{1/4}\FF_{2^d}$) empties the
cell (subspace base, 2-torsion quotient) on every prime-degree standardised
curve. The production 2-torsion census added no relations. A September
study measured the same bound from the other side: the invariants used by
translation-invariant bases on $y^2+xy=x^3+1$ are pullbacks by $1+\tau$,
$1-\tau$ and $[2]$, held on all 10{,}238 points at $n=7,11,13$, and each
kernel has at most 4 points, so they change a decomposition by at most 2
bits per summand (\art{research/cheap_closures_20261007/README.md}).

\section{Free relations}\label{sec:kc-free}

Pair-table self-collisions give relations with right-hand side zero: four
base points that sum to $O$, verified as exact group identities
(\art{RESEARCH_FREE_RELATIONS_CENSUS_20261012.md}). The relations exist and
are highly correlated (Table~\ref{tab:kc-free}): the 588 nontrivial rows at
$n=61$, $K=600$ touch 168 of 600 columns, and the 72 at $n=83$ span six
dimensions.

\begin{table}[t]
\centering\small
\begin{tabular}{@{}rrrrrr@{}}
\toprule
$n$ & $K$ & collisions & verified nontrivial & homogeneous rank & needed \\
\midrule
61 & 600 & 10{,}188 & 588 & 49 & 599 \\
61 & 1{,}200 & 12{,}732 & 4{,}800 & 746 & 1{,}199 \\
61 & 1{,}500 & 10{,}924 & 6{,}000 & 1{,}197 & 1{,}499 \\
83 & 600 & 9{,}672 & 72 & 6 & 599 \\
\bottomrule
\end{tabular}
\caption[]{Free relations from pair-table self-collisions: the rank they
reach against the $K-1$ the rank stage needs. Source:
\art{RESEARCH_FREE_RELATIONS_CENSUS_20261012.md}.}
\label{tab:kc-free}
\end{table}

\begin{itemize}
\item \emph{Rank law.} Rank grows about as $K^4$ with a small constant, so a
  free rank stage needs $K\approx1{,}500$--$1{,}900$ at $n=61$ and
  $K\approx2{,}600$ at $n=83$, where the census scan and the $K^2$ linear
  algebra already exceed the guided rank stage (13.1\,s at $n=61$, 4.8\,min
  on 12 threads at $n=83$), before 36\,GB of memory at $n=83$.
\item \emph{In the S3-root pipeline.} On by default for $n\le63$, the
  harvest costs $0.960\times$ $[0.946,0.976]$ of the tuned incumbent pooled
  over $n=41$, 53, 61, in S3 root solves per complete solve; wall time
  agrees, $0.941\times$
  (\art{research/ic_improve_20261015/02-freerel-production/README.md}).
\item \emph{In the pair-table pipeline.} Under rule v4 the harvest is
  correct (6{,}187 relations, none rejected, in 510 runs) but yields 2--11
  relations per solve against 58--327 needed, cuts trials 0.5--9\%, and makes
  the build $1.4$--$1.55\times$ slower: the cold cost rises by $5.3$--$8.8\%$
  at every size ($1.053$--$1.088$ over 72 pairs per size). The $0.960$ does
  not transfer, because that pipeline costs 59\,s per solve at $n=61$
  against 3.7\,s here (\art{research/ic_tool_program/rule/v4/README.md}).
\end{itemize}

\section{Batch index calculus and restart portfolios}\label{sec:kc-batch}

The many-target accounting correction and its batch-rho reference are in
\S\ref{sec:rv-batch}; Table~C is in \S\ref{sec:tablea}. The lead records
two statements that bound the route
(\art{crypto-autoresearcher:knowledge/leads/ecdlp/ic-binary/LEAD-ICB-multi-target-and-portfolio.yaml}):
\begin{itemize}
\item \emph{The generic multi-instance bound covers set-up.} $L$ logarithms
  cost at least about $\sqrt{Lr/\abs{\mathrm{Aut}}}$ for a generic-oracle
  index calculus including its set-up, and Kuhn--Struik batch rho attains
  it; the $L^{-1}$ amortisation applies only to the online part, after a
  set-up that obeys the same bound. This is the multi-target form of the
  Corrigan-Gibbs--Kogan bound~\cite{cgk2018} of \S\ref{sec:theory}.
\item \emph{A restart portfolio cannot beat pure rho} when the
  index-calculus cost concentrates at a floor ratio of at least 1. The
  ledger hypothesis \texttt{H-BINSTD-32867d} is marked supported as a scoped
  observational hold, strength preliminary (\texttt{EV-BINSTD-2d34cb}); its
  instrument was a deterministic closure, a controlled null rather than a
  solver.
\end{itemize}
The four-target union column of the September batch thread is a rare-event
linearisation, not a batch speed-up. Table~C's 32-target rows stay
$3.62$--$14.9\times$ batch rho, with the reference taken from a different
session (S20) at the canonical step, and all nine failed the 24 October
gate.

\chapter{The $n=83$ campaign after 21 October}\label{ch:n83late}

\S\ref{sec:n83} covers degree 83 to 21 October and the 24 October
field-table correction, after which every Weil-descent engine's earlier
$n=83$ refutation is void. Forty-one study directories were added from
24 October to 1 November. None of them admits a new IC/rho ratio, and none
is a leaderboard or \texttt{benchmarks/} row. They fall into eight groups
(Table~\ref{tab:n83l-groups}), which this chapter takes in four parts:
\begin{itemize}
\item factor-base selection, index capacity and controls on the pipeline
  curve (\S\ref{sec:n83l-pipeline});
\item decomposition at $n=83$ after the correction
  (\S\ref{sec:n83l-sat});
\item certified degree bounds for the descended systems
  (\S\ref{sec:n83l-degree});
\item the genus-3 image transfer (\S\ref{sec:n83l-genus3}).
\end{itemize}

\paragraph{Two curves, three families.} Degree 83 means two different
targets in these studies, and their numbers must not be mixed.
\begin{itemize}
\item \emph{The pipeline curve} \texttt{EC1N83Ckb1a1haacb69ee9f37},
  $y^2+xy=x^3+x^2+1$ ($a=1$), has a 53-bit subgroup of order
  $8{,}569{,}786{,}107{,}849{,}059$ and a cofactor of about
  $1.1\cdot10^9$. Every K-series base, every online and cold time and the
  matched precomputed-rho control are on this curve.
\item \emph{The primary curve} $y^2+xy=x^3+1$ ($a=0$, modulus
  $z^{83}+z^{45}+z^2+z+1$, trace $-6{,}151{,}469{,}093{,}347$) has the 81-bit
  subgroup $r=2{,}417{,}851{,}639{,}230{,}796{,}216{,}685{,}689$ and cofactor
  4. The genus-3 transfer and the sixteen factor-base designs target it.
\item \emph{The degree sweeps} use full-curve affine bases on three
  families at the same size: the Koblitz curve, a general binary curve and an
  83-bit prime curve, with no subgroup filter.
\end{itemize}

\begin{table}[t]
\centering\small
\begin{tabularx}{\linewidth}{@{}>{\raggedright}p{0.22\linewidth}r>{\raggedright}X>{\raggedright\arraybackslash}p{0.2\linewidth}@{}}
\toprule
Group & Dirs & What it established & Status \\
\midrule
Nearby grid K1102--K1120 & 4 & 19 exact prefix bases; 16-point confirmation admitted 112 of 112 IC/rho rows; no base selected & Tier~1 cloud; \textsc{bounded\_unknown} selection \\
Current-source K-series & 10 & K2304--K5500 bases reconstruct exactly; one K2304 run completes; K3596 and K4405 capped & local controls \\
Index capacity layouts & 7 & four-byte, split and fingerprint root indices; K5500 in 25.61--28.61\,GiB of allocations & allocations calculated, not measured \\
Cross-family bridge & 2 & 16 sealed points; never dispatched (Spot quota) & no timing \\
Precomputed-rho control & 1 & rho faster online in all three matched arms & Table~\ref{tab:distance-n83} \\
Decomposition engines & 7 & chained-$\Sem3$ SAT decides $m=2$--9 on a 167-point base with tables and a meet & stage diagnostic \\
Degree certificates & 7 & first fall 2; exact top-form degrees 28 and 111--609; no regularity computed & producer certificates \\
Genus-3 image & 3 & executable transfer; Prym coprime to $r$; degree-2 fibres carry no information & structural, negative \\
\bottomrule
\end{tabularx}
\caption[]{The 41 $n=83$ study directories of 24 October--1 November, by
group. ``Tier~1'' is the benchmark-contract host tier of
\art{docs/IC_BENCHMARK_TIERS.md}. Sources: the \texttt{research/n83\_*}
directories cited in this chapter.}
\label{tab:n83l-groups}
\end{table}

\section{Factor-base selection on the pipeline curve}\label{sec:n83l-pipeline}

\paragraph{The nearby grid.} The ordered K1120 producer dump fixes every
complete signed-Frobenius orbit prefix from K1102 to K1120; all 19 candidate
manifests were reconstructed from it, and K1104, K1107 and K1120 match their
published dumps byte for byte
(\art{research/n83_nearby_full_grid_20261024/README.md}). The same-point
Tier~1 screen of all 19 (84 children: 60 IC and 24 rho, each IC block
bracketed by K1107 and a rho control) failed its frozen 4\,GiB disk gate on
the first two hosted attempts, with 2.26 and 2.93\,GiB free and zero
measured children (\art{research/n83_nearby_grid_tier1_20261024/README.md});
an AWS Spot variant was blocked by the provider's Spot instance quota
(\texttt{Client.MaxSpotInstanceCountExceeded}, status
\texttt{BLOCKED\_PROVIDER\_SPOT\_QUOTA};
\art{research/n83_nearby_grid_aws_tier1_20261024/admission_20261024.json}).
A later hosted screen run (GitHub run 40150483792) supplied the
confirmation's candidate set.

\paragraph{The 16-point confirmation.} Sixteen fresh public points, sealed
before the screen result was inspected, were run on seven candidates in two
eight-point Tier~1 strata (AMD EPYC 7763 and Intel Xeon Platinum 8573C
runners), each IC run paired with a same-point rho
(\art{research/n83_nearby_grid_confirm16_20261024/RESULT.md}). The
independent verifier admitted 112 of 112 rows, and a local replay of the raw
artifacts reproduced the hosted aggregate byte for byte.
Table~\ref{tab:n83l-confirm} gives the panel. K1115 has the lowest
geometric-mean online time ($47.114$\,ms) and K1111 the lowest mean cold
process ($159.263$\,s), but the frozen sign-flip rule with Bonferroni
$\alpha=0.05/21$ returned \textsc{bounded\_unknown} for both metrics (K1115
against K1111 online, $p=0.959$; K1111 against K1116 cold, $p=0.773$). No
base was selected. The rho/IC column is an online reading with set-up
excluded and a table-less rho, the comparison \S\ref{sec:koblitz} retired;
on the same panel the cold IC process takes 159--166\,s, while the
same-point table-less rho online times at point 0 are
1{,}455--3{,}715\,ms.

\begin{table}[t]
\centering\small
\begin{tabular}{@{}rrrrr@{}}
\toprule
K & base points & online, ms (geo.\ mean) & cold, s (mean) & median rho/IC online \\
\midrule
1103 & 183{,}098 & 142.358 & 162.320 & 16.632 \\
1107 & 183{,}762 & 72.916 & 160.806 & 28.414 \\
1109 & 184{,}094 & 99.303 & 165.435 & 14.045 \\
1110 & 184{,}260 & 77.205 & 165.329 & 28.296 \\
1111 & 184{,}426 & 48.226 & 159.263 & 45.384 \\
1115 & 185{,}090 & 47.114 & 165.754 & 27.967 \\
1116 & 185{,}256 & 58.977 & 159.395 & 56.961 \\
\bottomrule
\end{tabular}
\caption[]{The nearby-grid confirmation on the pipeline curve: 16 fresh
public points, two Tier~1 strata, 112 admitted IC/rho rows. Base points are
$166K$ subgroup-usable points in $K$ folded columns. The rho/IC column is
online only, against a rho with no table, set-up excluded. Source:
\art{research/n83_nearby_grid_confirm16_20261024/RESULT.md}.}
\label{tab:n83l-confirm}
\end{table}

\paragraph{A local five-base screen.} On one new public point, five
finalists each ran two fresh cold processes in forward and reverse order
(\art{research/n83_local_five_base_screen_20261024/RESULT.md}). Median cold
times were 66.2\,s (K992/root11, 164{,}672 points), 62.7\,s (K1024), 65.8\,s
(K1107), 119.7\,s (K2304) and 383.6\,s (K3300/wide3), with pair spreads of
4.7--19.1\%; the two same-point rho controls took 2.015 and 0.989\,s with
different seeds. Every hash, base, rank, relation and scalar replayed, and
no winner was selected because the K1024--K1107 gap (3.1\,s) is smaller than
the within-arm spread.

\paragraph{The K-series, re-identified.} The large bases of 14--21 October
were reconstructed from current-source parent dumps: the first 2{,}304,
3{,}300, 3{,}450 and 3{,}500 orbits of the K3596 dump (596{,}936 points)
and the first 4{,}300 and 4{,}405 of the K5500 dump (913{,}000 points)
reproduce the historical K2304 (382{,}464 points), K3300 wide3 (547{,}800),
K3450 wide3 (572{,}700), K3500 root4 (581{,}000), K4300 (713{,}800) and
K4405 (731{,}230) bases exactly, by ordered-dump and point-set SHA-256 and
the producer's BLAKE3 base hash
(\art{research/n83_k2304_current_source_prefix_20261025/},
\art{research/n83_k3300_current_source_prefix_20261025/},
\art{research/n83_crossfamily_current_20261025/README.md},
\art{research/n83_k4300_k4405_current_source_prefix_20261025/}). Full-process
controls on an earlier public point then gave
(Table~\ref{tab:n83l-controls}) one completed large base and three that did
not finish on the hosts available.

\begin{table}[t]
\centering\small
\begin{tabularx}{\linewidth}{@{}l>{\raggedright}Xrr>{\raggedright\arraybackslash}p{0.22\linewidth}@{}}
\toprule
Base & Host and cap & Wall, s & Peak RSS & Outcome \\
\midrule
K2304 root11 & local, 1{,}200\,s, 12\,GiB & 438.5 & 5.93\,GB & completed; independent replay passed; target stage 420.2\,ms \\
K1107 root4 & local macOS ARM64, 1{,}200\,s & 438.9 & 9.94\,GB & completed; replay passed; online 46.3\,ms \\
K3596 root4 & busy 48\,GiB macOS, 1{,}200\,s & 792.5; 1{,}201.0 & 0.80; 9.58\,GB & SIGTERM after CPU starvation; time cap \\
K4405 wide3 & local, 1{,}200\,s & 873.2 & 11.68\,GB & stopped by hand under host swap \\
K4405 wide3 & Linux Pod, 64\,GiB cgroup & --- & --- & not started; approved GPU pair out of capacity \\
\bottomrule
\end{tabularx}
\caption[]{Full-process controls on the pipeline curve, one earlier public
point each (fixture zero), outside the sealed bridge points; correctness and
resource controls, not Tier~1 pairs. K1107 root4 uses a fixed 8\,GiB root
table for 101{,}708{,}823 distinct roots. Sources:
\art{research/n83_k2304_current_source_control_20261026/README.md},
\art{research/n83_root4_small_control_20261025/README.md},
\art{research/n83_k3596_current_source_full_20261025/README.md},
\art{research/n83_k4405_current_source_full_control_20261026/README.md},
\art{research/n83_k4405_linux_control_20261026/README.md}.}
\label{tab:n83l-controls}
\end{table}

\paragraph{Index capacity.} Seven directories made the root index fit
larger bases. Each is an opt-in layout with an exact-key replay rule (a
candidate slot reconstructs and compares the full 83-bit root before a hit
is accepted), and each stops at calculated allocations for K5500:
\begin{itemize}
\item \emph{Exact four-byte slots} (31-bit state index and a root selector
  in one \texttt{u32}): principal allocations of $32{,}062{,}766{,}258$
  bytes, 29.86\,GiB, for K5500
  (\art{research/n83_k5500_fourbyte_20261025/README.md}).
\item \emph{ID31 four-byte slots} at the K5500/K4549 split: 12\,GiB of
  tables instead of 15, and state, table and Bloom arrays of 27.86\,GiB
  against 30.86\,GiB for five-byte ID31
  (\art{research/n83_fourbyte_id31_20261027/README.md}).
\item \emph{A relative-ID split} (primary $2^{31}$ and extension $2^{30}$
  slots, 1{,}255{,}377{,}750 eleven-byte states): 25.61\,GiB of principal
  allocations, 4.25\,GiB below monolithic ID31
  (\art{research/n83_k5500_split4_20261025/README.md}).
\item \emph{Fingerprint sidecars.} A two-bit sidecar brings the ID31 layout
  to 28.61\,GiB, still 2.25\,GiB below five-byte ID31, with the packing and
  capacity arithmetic proved in Lean~4.34.1
  (\art{research/n83_fourbyte_fp2_20261027/README.md}). A one-byte sidecar
  kept every digest identical and built the index faster in ABBA order on a
  shared ARM64 host: geometric-mean baseline over filtered ratios of
  $1.831$ (K128), $1.347$ (K512) and $1.505$ (K1024), with a K2048 baseline
  that drifted 42\% between bookends
  (\art{research/n83_k5500_fp8_20261025/README.md}).
\end{itemize}
Two whole-process ABBA controls on the K1107 base read the sidecars:
the one-byte sidecar gave baseline over fingerprint $1.677$ for the full
process and $2.157$ for one-target online, with a 15.3\% baseline drift
between bookends (\art{research/n83_fp8_e2e_control_20261025/README.md});
the two-bit sidecar gave whole-process ratios of $1.017$ and $1.294$ in the
two orders on a host whose load average ran from 22 to 276
(\art{research/n83_fourbyte_pair_20261027/README.md}). Both are local
diagnostics that the Tier~1 gate would not admit.

\paragraph{The bridge that did not run.} The cross-family bridge sealed 16
new public points for a same-point Tier~1 comparison of the small-base and
wide-base selections, with its nomination rule frozen before either panel
reported (\art{research/n83_crossfamily_bridge_20261025/README.md}). Its
current-source panel (K3450 wide3 and K3500 root4, two \texttt{g5g.4xlarge}
Spot instances, a \$4.20 cap) moved from \texttt{us-east-1d}, which failed
its placement-score gate, to \texttt{us-west-2c}; the README still records
that the panel has not dispatched or selected a winner
(\art{research/n83_crossfamily_current_20261025/README.md},
\art{research/n83_crossfamily_current_20261025/AWS_WEST2_ADDENDUM.md}).

\paragraph{The control that did run.} The matched precomputed-rho
reproduction of 28 October on K1107/root11 found rho faster online in all
three matched arms (Table~\ref{tab:distance-n83}), every interval above 1,
over 16 targets with every scalar replayed
(\art{research/n83_k1107_bl_reproduction_20261028/RESULT.md}). Its plain
table-less rho, a positive control only, reads rho/IC $10.7$
$[6.33,88.8]$ online, the same kind of reading as the rho/IC column of
Table~\ref{tab:n83l-confirm}.

\section{Decomposition at $n=83$ after the correction}\label{sec:n83l-sat}
The engines run at $n=83$ are described in \S\ref{sec:bo-n83}; this section records the decomposition results after the field-table correction.


After the 24 October correction (\S\ref{sec:n83}) chained-$\Sem3$
full-width SAT was the only decomposer that ran at $n=83$, and it decided
only $m=2$. Six preregistered follow-ups of 31 October--1 November extended
it on the pinned pipeline-curve fixture, with the factor base the Frobenius
union of $\mathrm{span}(1,z)$, 167 points
(\art{research/n83_chain_sat_meet_20261101/PROTOCOL.md}). Each cell has
three public targets (not expected to decompose) and two or five planted
ones. Table~\ref{tab:n83l-sat} gives the sequence. Every run had zero
soundness flags: no planted target was refuted and every witness re-adds in
the group.

\begin{table}[t]
\centering\small
\begin{tabularx}{\linewidth}{@{}l>{\raggedright}X>{\raggedright\arraybackslash}X@{}}
\toprule
Study & Change & Outcome \\
\midrule
order (31 Oct) & all-summand ordering; root-finder tables & $m=2$: conflicts halved (median public 13{,}461 $\to$ 7{,}195; 12.8 $\to$ 6.2\,s); $m=3$: tables refute every public target with 0 conflicts and find both planted; $m=4$: public refuted, no planted witness in 120\,s; $m\ge5$: capped \\
meet (1 Nov) & a meet-in-the-middle lookup for the middle intermediate & $m=5$--8: public 3/3 refuted with 0 conflicts (median 0.80--76.7\,s), planted 0/2; $m=9$: capped; builds 0.68\,s ($m=5,6$), 26.1--27.2\,s ($m=7,8$), 958.8\,s ($m=9$) \\
rebuild (1 Nov) & rebuild the planted witness from the lookup hit & $m=5$--8: planted 5/5 with 0 solver calls (median 0.91--81.5\,s) \\
indexed (1 Nov) & indexed lookup, rebuild by peeling & $m=8$: 78.4 $\to$ 39.7\,s; $m=9$: from cap to 3/3 and 5/5 (54.1\,s), depth-5 set of 372{,}406{,}144 fingerprints in 999\,s \\
parallel (1 Nov) & 4-thread back enumeration & $m=8$: 41.7 $\to$ 11.4\,s; $m=9$: 57.0 $\to$ 13.7\,s; verdicts identical \\
parbuild (1 Nov) & 4-thread depth-5 build & 1{,}006.8 and 993.4\,s $\to$ 287.7 and 283.3\,s, $3.50\times$ \\
\bottomrule
\end{tabularx}
\caption[]{Chained-$\Sem3$ SAT on the pipeline curve, $m=2$--9, 167-point
base; decomposition stage only, L0 four-vCPU hosts. Sources:
\art{research/n83_chain_sat_order_20261031/README.md},
\art{research/n83_chain_sat_meet_20261101/README.md},
\art{research/n83_chain_sat_rebuild_20261101/README.md},
\art{research/n83_chain_sat_indexed_20261101/README.md},
\art{research/n83_chain_sat_parallel_20261101/README.md},
\art{research/n83_chain_sat_parbuild_20261101/README.md}.}
\label{tab:n83l-sat}
\end{table}

The sequence turns SAT into a meet-in-the-middle with a SAT residue, and
then removes the residue: once the lookup decides the middle intermediate,
the public targets are refuted with zero conflicts because the lookup domain
is empty on every one of them and non-empty (10--250 entries) on every
planted one. The meet study states the cost as about $\abs{B}^{m-d}$ per
query after a $\abs{B}^d$ build. That is the generic floor, and the
improvements in Table~\ref{tab:n83l-sat} are constants of the enumeration
(ordering, tables, peeling, threads). The matched rho on this curve is
1.51\,s mean online. With 167 points and a 53-bit subgroup, a decomposition
over this base is rare, so these are stage measurements: no relation
collection, no complete logarithm and no rho ratio.

\section{Certified degree bounds}\label{sec:n83l-degree}

The descended $n=83$ systems were characterised exactly, not solved. Seven
directories, all producer certificates with fresh-process replay and
independent scientific review pending, give the following.

\paragraph{Solvers do not finish.} The balanced degree sweep (Koblitz,
general binary and prime; $m=2$--9; candidate domain $2^{\lfloor83/m\rfloor}$;
a planted and a natural target per cell, two computations each) ran 96
Singular attempts with 40\,s and 3\,GiB per solver: 53 hit the time limit,
43 the memory limit, and no top-form or affine basis degree completed
(\art{research/n83_ic_degree_sweep_20261027/README.md}). A fixed-base rerun
at $\ell=27$, on full-curve bases that enumerate all $2^{27}=134{,}217{,}728$
candidate $x$-values per curve
(\art{research/n83_fixed_factor_base_20261028/README.md}), completed 4
affine bases (all $m=2$, Koblitz and binary), with 63 time and 29 memory
exits (\art{research/n83_fixed_degree_sweep_20261028/README.md}).

\paragraph{First fall is 2 on binary curves.} On every retained input the
Koblitz and general binary systems have certified first fall degree 2 at
every $m=2$--9, planted and natural (32 exact witnesses; all 48 replays
passed), consistent with Corollary~\ref{cor:ffd2}. On
the prime curve the balanced presentation has the proof-assisted lower
bounds 24, 17, 14, 12, 10, 9, 9, 8 for $m=2$--9
(\art{research/n83_ic_degree_sweep_20261027/README.md}). The fixed
$\ell=27$ catalogue certifies 48 first falls and 32 degree-two closures: at
$m=2$ the binary systems have first and last fall 2 (a restricted global
certificate; in each planted case $V_2$ has 53 independent linear
constraints on 54 variables), at $m=3$--9 the last fall is unresolved with
lower bound 2, and the prime systems have first fall at least 5
(\art{research/n83_fixed_fall_catalog_20261028/README.md}). A later
degree-four witness census raised the certified last-fall lower bound of
the twelve retained $m=3$ Koblitz and binary systems from 3 to 4, with 249
strict degree-four coordinate witnesses each
(\art{research/n83_candidate_pdp_20261028/README.md}).

\paragraph{Top-form degrees are large and exact.} The top-form graded
quotient was bounded for all 48 fixed-base systems
(\art{research/n83_top_quotient_20261028/README.md}) and the binary chain
bounds were then closed exactly by a written overlap proof with native
checks (\art{research/n83_candidate_pdp_20261028/README.md}).
Table~\ref{tab:n83l-top} gives both. For the prime systems at $m\ge3$ the
top quotient is positive-dimensional, with no finite first vanishing degree
in the saved presentation; the source says this is not a claim of infinite
regularity or about the affine ideal. At $m=2$ the binary systems have
$\text{FFD}=\text{LFD}=2$ and top-form degree 28, which separates those
quantities explicitly for these presentations.

\begin{table}[t]
\centering\small
\begin{tabular}{@{}rlr@{}}
\toprule
$m$ & certified bounds, Koblitz and binary & exact top-form degree \\
\midrule
2 & $28\le d_{\mathrm{reg}}\le55$ & 28 (prime: 42) \\
3 & $82\le d_{\mathrm{reg}}\le165$ & 111 \\
4 & $109\le d_{\mathrm{reg}}\le275$ & 194 \\
5 & $136\le d_{\mathrm{reg}}\le385$ & 277 \\
6 & $163\le d_{\mathrm{reg}}\le495$ & 360 \\
7 & $190\le d_{\mathrm{reg}}\le605$ & 443 \\
8 & $217\le d_{\mathrm{reg}}\le715$ & 526 \\
9 & $244\le d_{\mathrm{reg}}\le825$ & 609 \\
\bottomrule
\end{tabular}
\caption[]{Top-form degrees of the fixed $\ell=27$ descended systems at
$n=83$, both targets. Bounds: certificates of 28 October; exact values:
the overlap proof of 29 October, which closes every interval at
$111+83(m-3)$ for $m\ge3$. Neither is an affine solving degree. Sources:
\art{research/n83_top_quotient_20261028/README.md},
\art{research/n83_candidate_pdp_20261028/README.md}.}
\label{tab:n83l-top}
\end{table}

\paragraph{Encodings and candidate bases.} For the prime candidate
interval a radix census over all 3{,}010 partitions of 27 (35 exact Hilbert
records and 2{,}975 certified lower bounds) found the unique minimum exact
first vanishing degree, 42, in the all-binary presentation, which excludes
improvements within that presentation family
(\art{research/n83_candidate_pdp_20261028/README.md}; the original radix
protocol's unequal-arm accounting mismatch is recorded in a correction file).
The same directory certifies $7\le\text{FFD}\le17$, conditionally, for 42
prime $m=3$--9 systems. A separate design study lists sixteen factor-base
constructions for the primary $a=0$ curve with a preregistered protocol and
proof obligations, and no measurement
(\art{research/n83_factor_base_designs_20261031/README.md}).

\paragraph{Reading.} The binary systems fall at degree 2 and their top-form
degree grows by 83 per summand. These are certified properties of
specific presentations, not attack costs. They say nothing for or against a
low affine solving degree, but they show that the descended $n=83$ systems
are far outside the reach of the solvers that were run, which matches the
96 and 92 unfinished attempts above.

\section{The genus-3 image transfer}\label{sec:n83l-genus3}

A same-field degree-3 cover $\phi: C\to E$ with $C$ of genus 3 descends to
$\FF_2$: $E: y^2+xy=x^3+1$ and $C: v^2+u^2v=u^7+u$, with equal point counts
over $\FF_2$, $\FF_4$ and $\FF_8$
(\art{research/n83_jacobian_projector_20261031/RESULTS.md}). Three studies
of 31 October--1 November asked whether pulling the $n=83$ instance into
$\mathrm{Jac}(C)$ gives anything rho does not have.
\begin{itemize}
\item \emph{The transfer executes}
  (\art{research/n83_genus3_image_solver_20261031/README.md}). The exact
  $n=83$ pullbacks of $P$ and $Q=[37]P$ are reduced Mumford classes of exact
  prime order $r$; typed certificates prove $\phi_*\phi^*=[3]$, so the
  pullback is injective on the subgroup. Rho and BSGS in the Jacobian image
  solved toy instances at $n=9,13,19,23$ without the planted scalar, all
  400 destination attempts verified ($n=7$ failed and is retained). The
  source and destination operation counts have no certified common unit, so
  no advantage is promoted; the full-Jacobian genus-3 index calculus does
  not apply to an image of order $r\approx q$, and its ambient $q^{4/3}$
  ledger loses to source rho by about 70.17 exponent bits.
\item \emph{The Prym shares nothing with the target}
  (\art{research/n83_jacobian_projector_20261031/RESULTS.md}). With the
  closed-point norm, $e=\phi^*\phi_*$ satisfies $e^2=[3]e$, and on
  $r$-torsion $p=[3^{-1}\bmod r]e$ is a checked idempotent. At $n=83$ the
  Prym has Frobenius polynomial
  $X^4+q^2=(X^2-2^{42}X+q)(X^2+2^{42}X+q)$, with elliptic factors of traces
  $\pm2^{42}$, neither equal to $E$'s trace, and
  $\gcd(r,q+1\mp2^{42})=\gcd(r,q^2+1)=1$. There is no $r$-primary route
  through the two-dimensional complement. The identities are checked in
  Lean.
\item \emph{Degree-2 fibres are tautological}
  (\art{research/n83_image_closed_degree2_20261101/README.md}). On every odd
  extension tested a generic rational fibre has closed points of degrees
  $1+2$ whose elliptic labels are exactly $P$ and $[2]P$, so the projected
  identity is $[3]P=P+[2]P$ and carries no logarithmic information. The
  frozen information gate fails, so no collector or timing run is admitted;
  ordinary ambient genus-3 collection and fixed bounded-degree branching
  also fail the comparison with source rho. A growing-degree or adaptive
  target-local router is left open.
\end{itemize}
This is the binary counterpart of the cover results of
\S\ref{sec:transfer}: the map is real and certified, and the destination
carries the same $r$-part and nothing that makes it cheaper.

\chapter{Weil descent and degree, continued}\label{ch:binarydegree}

\S\ref{sec:binary-descent} measures the degree of Weil-descended summation
systems to 24 October: first and last fall degrees stay at 3--4, the
refutation degree tracks $\ell$ within one of semi-regular, closure
algorithms refute a degree or two earlier but rise past 4 at $m=3$,
$\ell=7$, and guessing one or two variables restores degree 4 there.
\S\ref{sec:frontier} adds the first $\ell=8$ runs: four guessed variables
bring all four unsatisfiable $(22,8)$ draws back to degree 4, one
three-variable set fails on three of four, and the hybrid beats degree 5 at
the priced ECC2K-130 cell only if at most six guesses suffice. This chapter
records what followed:
\begin{itemize}
\item which guessed coordinate matters at $\ell=8$, not only how many
  (\S\ref{sec:bd-guess});
\item the first measured degree-5 closure cell (\S\ref{sec:bd-degree5});
\item the per-attempt criterion and the companion repository's closure and
  SAT pilots (\S\ref{sec:bd-criterion});
\item a direct check of Semaev's 2015 first-fall statement on his own cells
  (\S\ref{sec:bd-semaev});
\item the ceiling on learned SAT guidance (\S\ref{sec:bd-sat});
\item the finding that F6-IC with a pair index is the meet-in-the-middle
  scan (\S\ref{sec:bd-f6}).
\end{itemize}
Every result here is exact algebra or a counted stage diagnostic on toy
fields, run on unisolated 4-core, 16\,GB containers (the RunPod degree-5
cell excepted). None has a validation file. One carries a proof-carrying
claim, recorded as unverified. Closure outcomes do not depend on the host;
timings are indicative only.

\section{Guessing at $\ell=8$: which coordinate, not how many}\label{sec:bd-guess}
The closure and guessing instruments are described in \S\ref{sec:in-closure}; this section records what they measured at $\ell=8$.


All studies in this section use the frozen $m=3$, $(n,\ell)=(22,8)$ draws
of 22 October ($N=46$ unknowns, 44 equations, surplus $-2$; unsatisfiable
draws 0, 3, 5, 6) or fresh draws from the same generator, the MutantXL-style
closure \texttt{dreg\_closure} or its guess variants, and a registered
reading of each branch as refuting at degree 4 (the constant 1 appears) or
not (degree 4 saturates without it). Table~\ref{tab:bd-guess} lists them.

\begin{table}[t]
\centering\small
\begin{tabularx}{\linewidth}{@{}l>{\raggedright}X>{\raggedright}X>{\raggedright\arraybackslash}p{0.2\linewidth}@{}}
\toprule
Study & Guess set (all set to 0) & Outcome & Verdict \\
\midrule
guess (22 Oct) & $x_7,x_{15},x_{23},x_{24}$: one top coordinate per summand and one intermediate & 64 of 64 branches refute at 4; suffices on 4 of 4 draws & four suffice (an upper bound) \\
guess3 (23 Oct) & $x_7,x_{15},x_{23}$ & draws 0, 3, 5 fail at branch 0 (closure rank 129{,}221, 128{,}213, 129{,}893 of 136{,}698); draw 6 suffices & three do not suffice for this set \\
null (28 Oct) & none, at $(22,7)$, $N=43$ & 4 of 4 draws do not refute at 4 & 12 of 12 unspecialised $\ell=7$ draws now fail \\
invariant (29 Oct) & S3: $a_7+b_7$, $c_7$, $e_1[0]$; X3: $a_7$, $c_7$, $e_1[0]$ & S3 suffices on 2 of 4, X3 on 1 of 4; S3 never worse & mixed \\
invariant, more (29 Oct) & S3 against X3 on 8 fresh draws & S3 5 of 8, X3 0 of 8; 5 of 5 discordant draws favour S3, $p=0.031$ & S3 better (preregistered) \\
pair (29--30 Oct) & D3: $a_7+b_7$, $a_6+b_6$, $e_1[0]$ against S3 & D3 2 of 12, S3 7 of 12; 5 discordant, all S3; two-sided $p=0.0625$ & no detected difference \\
second link (5 Nov) & L3: $a_7+b_7$, $(e_1+x_3)_{20}$, $e_1[0]$ against S3 & L3 6 of 12, S3 7 of 12; one discordant draw; $p=1.0$ & no detected difference; unverified claim \texttt{PCC-\allowbreak20261106-\allowbreak607061} \\
\bottomrule
\end{tabularx}
\caption[]{Guessed-variable closures at $m=3$, $(n,\ell)=(22,8)$ (the null
row at $(22,7)$), degree 4, unsatisfiable draws only; $a,b,c$ are the
summands' coordinates and $s=a+b$ the swap invariant of the first link.
``Suffices'' means every branch of the draw refutes at degree 4. Sources:
\art{research/dreg_closure_guess_20261022/RESULTS.md},
\art{research/dreg_closure_guess3_20261023/RESULTS.md},
\art{research/dreg_closure_null22_7_20261028/RESULTS.md},
\art{research/dreg_invariant_guess_20261029/RESULTS.md},
\art{research/dreg_invariant_guess_more_20261029/RESULTS.md},
\art{research/dreg_invariant_pair_guess_20261029/RESULTS.md},
\art{research/dreg_second_link_guess_20261105/RESULTS.md}.}
\label{tab:bd-guess}
\end{table}

\paragraph{What the series shows.} At equal guess count, guessing a
coordinate of the swap invariant $s=x_1+x_2$ refutes at degree 4 on more
draws than guessing a raw summand coordinate: preregistered on fresh draws,
$p=0.031$ one-sided, and 6 of 6 discordant draws over both studies
($p=0.016$, secondary). S3 suffices on 7 of 12 draws in all and X3 on 1 of
12. This contradicts the summand-dimension reading proposed after the
three-guess study: S3 leaves both summands at dimension 8 and suffices more
often than X3, which lowers $x_1$ to dimension 7. Putting a second
invariant guess on the same link (D3) is never better than the raw guess on
$x_3$; putting it on the second link's invariant $e_1+x_3$ (L3) is as good
as the raw guess on these 12 draws. The D3 failures all sit at branches 0
and 4, where both top coordinates of $s$ vanish. Four guesses
$\{7,15,23,24\}$ still suffice on every draw measured, and S3 fails on 5 of
12.

\paragraph{What it does not show.} Every row is one cell, $m=3$, with one
guess set per arm. The unspecialised $\ell=8$ system is not measured at
degree 4 (179{,}447 columns); ``does not refute'' means a closure degree of at
least 5, and degree 5 is not measured here. The null study's 12 of 12
unspecialised $\ell=7$ failures over $N=40$, 42, 43 and surplus $-2$, 0,
$+1$ remove the one confound the earlier $\ell=7$ registration kept (a
non-positive surplus), but cannot distinguish ``3 of 4'' from ``4 of 4''
against the three-guess $\ell=8$ result. The second-link study's own reading
is that one invariant guess per link extends along the chain, so $m=4$ is
its natural test; $m=4$ at $\ell=8$ has not run. The preregistered
high-arity discriminator that would run it
(\art{research/dreg_high_arity_20261101/README.md}) went through four
protocol versions, each superseded before data by a source or manifest
change on \texttt{main}; v4 is frozen and its launch is blocked.

\paragraph{The price, unchanged.} The guess-then-degree-4 oracle costs
$2^k\binom{N-k}{\le4}$ per target against $\binom{N}{\le5}$ for degree 5.
At the priced ECC2K-130 cell ($m=4$, $\ell=33$, $N=2\cdot131+4\cdot33=394$,
$w=1$) the ratio $\binom{394}{\le5}/(2^k\binom{394-k}{\le4})$ is $78.2$ at
$k=0$, $5.09$ at $k=4$, $1.30$ at $k=6$ and $0.66$ at $k=7$, recomputed in
the three-guess study. Both registered growth rules give $k\approx105$--108
there. Nothing in the series measures $k$ at $m=4$, so the $+8.43$ to
$+58.27$-bit range of \S\ref{sec:frontier} stands.

\paragraph{Other degree levers of the same weeks.} A preregistered
invariant arm that adds the summand-swap invariant $p$ as a variable makes
every equation quadratic but does not refute at degree 3 at $(22,6)$, where
the plain system refutes at 4; its degree-4 matrix would have $5.1$--$5.5\times$
the plain system's columns ($523{,}686$--$637{,}393$ against $102{,}091$), and at
$(21,7)$ degree 4 is unmeasured
(\art{research/dreg_pair_invariant_20261028/RESULTS.md}). Fixing one variable
usually brings the $\ell=7$ draws back to degree 4, but the null control
refuted 6 of 8 and was invalid, so the study is uninformative for its
registered question (\art{research/dreg_closure_specialise_20261017/RESULTS.md}).

\section{A measured degree-5 closure}\label{sec:bd-degree5}

The descent rows of Table~\ref{tab:prog_koblitz-ecc2k130} price the $m=4$
ECC2K-130 route at degree 4 ($+7.62$ bits, contradicted) and degree 5
($+10.32$ bits at $w=1$, $+45.71$ at $w=2$). Until November no degree-5
closure had been computed. The leads study of 24 October
(\art{research/koblitz_ic_leads_build_20261024/README.md}; Amendment~1 post
hoc, Amendments 2--7 frozen before the runs they govern) built four leads
in Rust and, in Amendment~7 (4 November), measured cell C1.
\begin{itemize}
\item \emph{C1 refutes at degree 5.} On $m=3$, $(n,\ell)=(19,7)$, draws 0
  and 3, the incremental engine \texttt{pe\_closure} (partial, low-first,
  batch 64) found 1 in the degree-5 closure, so each draw's closure degree is
  5. Draw 0: 607{,}530 of 760{,}099 columns filled at refutation (0.80),
  31{,}632\,s, 52.6\,GB peak resident; draw 3: 606{,}973 columns, 31{,}888\,s,
  52.5\,GB. Host: 32 vCPU AMD EPYC 9965 with 128\,GB (RunPod).
\item \emph{Memory.} The dense engine aborted on the same draws at 424 and
  423\,GB resident, more than twice its 217\,GB projection and about $8\times$
  the incremental engine's measured peak, which was below the derived
  73\,GB bound.
\item \emph{Scope.} One cell, $m=3$ not $m=4$, two of four draws, no
  satisfiable control (the owner chose draws 0 and 3). The study reads it as
  support for the degree-5 rows of
  \art{research/notes/ecc2k130/RESEARCH_ECC2K130_DESCENT_DEGREE.md} at
  $\ell=7$, and it changes no cost row. C1 cost about \$53 of an \$80 cap
  over six pods, of which only the last measured.
\end{itemize}
The other leads of the same study change no cost either. Partial
enlargement (L2) does not reduce memory in general: in the registered
order the basis fills 0.66--0.997 of the column space before 1 appears, and
other queue orders and batch sizes keep every run at 0.61--1.00. MXL's
variable restriction is the one exception, on $m=4$, $(11,5)$: all four
draws refute at fill 0.38--0.45 with $1.8$--$2.5\times$ fewer products, none
reaching the 0.25 payoff; elsewhere it does not help. The useful by-product
is the incremental engine itself, which holds one row per column plus one
chunk: 1.3--2.8\,GB resident at degree 4 against 1.9--4.2\,GB for the full
order.

\section{The per-attempt criterion and the closure pilots}\label{sec:bd-criterion}

The companion repository's degree pilots of 9--14 October are labelled by
their authors exploratory, ``not program evidence'', with no run record or
ledger write. They are worth recording for the criterion they state and for
three measurements that \S\ref{sec:binary-descent} does not have.

\paragraph{The criterion.} If a failed decomposition attempt costs
$2^{\alpha\ell}$, relation collection can beat rho only if $\alpha<m/2-1$
(\art{crypto-autoresearcher:analysis/ecdlp2m-iteration-20261009/m3-closure-pilot/README.md}).
A harvest of 411 measured rows and 1{,}554 per-target values from the
repository found that SAT and enumeration solvers spend about $2^{m\ell}$
per failed attempt
(\art{crypto-autoresearcher:analysis/ecdlp2m-iteration-20261009/ctrial-harvest/README.md}):
WDSat's unsatisfiable conflict count equals $2^{3\ell}/3!$ to within
0.03--1\% for $\ell=5$--11, slope $2.9985$ bits per unit $\ell$
$[2.9979,2.9992]$. At $m=3$ the budgets need $-1.7$ to $-3.1$ bits per unit
$\ell$ against a measured $+3.36$; at $m=4$, $-0.12$ to $-0.61$ against
$+2.32$ for meet-in-the-middle. Extrapolated to $n=131$ at the budget
$\ell$, WDSat with symmetry breaking lies at least 91 bits outside every
budget line, a Riemann--Roch $L(4O)$ oracle measured at $n=131$ at least
55, and the $t=4$ chained meet-in-the-middle at least 52; $m=5$, 6 had no
usable cells. Only a bounded-degree algebraic refutation could meet the
criterion.

\paragraph{The $m=4$ SAT exponent.} Run on the published
Trimoska et al.~\cite{trimoska2020} object and configuration, the pilot
reproduces $\alpha=2.981$ $[2.973,2.997]$ at $m=3$; on chained $\Sem3$ at
$m=4$ with Gaussian elimination it measures $\alpha=4.328$
$[4.315,4.340]$ ($4.310$ at fixed $n=17$), against the $\alpha<1$ the $m=4$
budget needs
(\art{crypto-autoresearcher:analysis/ecdlp2m-iteration-20261009/m4-sat-cost-pilot/README.md}).
The pilot also found that WDSat's stock \texttt{-b} mode returns ``SAT'' with
undefined bits on unsatisfiable $m=4$ instances.

\paragraph{Closure degree stays at 4 on toy chains, and is local.}
\begin{itemize}
\item \emph{$m=3$.} The degree-4 mutant closure $W_4$ refutes 60 of 60
  unsatisfiable chained-$\Sem3$ instances at $n=9$ and 40 of 40 at $n=11$,
  with no false refutation of a satisfiable one, where plain Macaulay $M_4$
  and $M_5$ refute none (plain needs degree 6 at $n=9$; semi-regular
  predicts 7--11). A support-matched null is refuted and a dense null is
  not, so the low-degree refutation follows the sparse, block-bilinear
  support of the descended system, not Semaev's coefficients.
\item \emph{$m=4$, 5.} $W_4$ refutes 25 of 25 at $m=4$, $n=7$; 10 of 10 at
  $n=9$; 2 of 2 at $n=11$; 10 of 10 at $m=5$, $n=7$; and $W_3$ refutes 30
  of 30 at $m=5$, $n=9$, $k=2$, while plain $M_4$ refutes nothing
  (\art{crypto-autoresearcher:analysis/ecdlp2m-iteration-20261009/m45-closure-pilot/README.md}).
\item \emph{Locality.} Restricting monomials to consecutive blocks
  (locality 1) keeps the refutation on 124 of 135 unsatisfiable instances
  with 39--69\% fewer columns at $m=3,4$; every miss refutes at locality 2,
  and none needed locality 3. At $m=5$, $n=9$, $k=3$ locality 1 refutes none
  of 10 and locality 2 all 5 tried
  (\art{crypto-autoresearcher:analysis/ecdlp2m-iteration-20261009/block-local-closure-pilot/README.md}).
\end{itemize}
These are the toy-scale version of the MutantXL result of
\S\ref{sec:binary-descent}, and they meet the same wall: the closure stays at
degree 4 while the columns stay within a host's memory, and the
$\ell=7$--8 cells of \S\ref{sec:bd-guess} and the degree-5 cell above are
where it leaves degree 4.

\section{Semaev 2015, \S4.5, on its own cells}\label{sec:bd-semaev}

Chapter~\ref{ch:falldegree} proves that the binary $m=2$ descent and
Semaev's chained system have first fall degree 2 for every $n$
(Corollary~\ref{cor:ffd2}) and notes that Semaev states 4 in his
\S4.5~\cite{semaev2015}. A frozen check of 2 November ran his own $t=2$
cells of Table~2, $S_3(x_1,x_2,z)$ with $x_1,x_2\in V$ and random $z$ and
$B$, 20 instances per cell and field polynomial, 200 in all
(\art{research/semaev_t2_trace_20261102/REPORT.md}; exact
$\FF_{2^n}$ and Boolean-ring arithmetic in standard-library Python; source
arXiv:1504.01175v1).
\begin{itemize}
\item \emph{Refuted: ``the first fall degree \dots\ is always 4''.} On all 200
  instances the first fall degree is 2, each with an explicit trace
  certificate: with $\lambda=z^{-2}$,
  $\operatorname{Tr}(z^{-2}S_3)=\operatorname{Tr}(x_1)+\operatorname{Tr}(x_2)+\operatorname{Tr}(Bz^{-2})$,
  a constant combination of the coordinate equations of degree at most 1.
\item \emph{Split: the corpus inference that these rows refute
  $d_{F4}\le d_{\mathrm{ff}}$.} The least degree at which the Boolean
  Macaulay span contains 1 is 2 on 154 instances and 3 on 46, never 4
  (per cell: $(15,4)$ 18/2 and 19/1; $(15,5)$ 20/0 twice; $(16,4)$ 18/2 and
  20/0; $(19,3)$ 6/14 and 4/16; $(21,3)$ 15/5 and 14/6). So the sharp
  first-fall assumption fails on 46 of 200 instances, not everywhere, and
  Semaev's reported $d_{F4}=4$ is Magma's step degree, not what the systems
  require.
\item \emph{Why.} The quadratic parts depend only on $x_1x_2\in V\cdot V$, of
  dimension $2k-1<n$, so most $t=2$ systems are inconsistent at degree 2 or
  3. None of the 200 has a solution in $V\times V$, which matches Semaev's
  success probabilities of 0.0006--0.0155 for these rows.
\end{itemize}
It complements the Semaev 2015 audits of Chapter~\ref{ch:summer}, which
tested eq.~(11) and the typed-image ratio but not the first-fall statement:
on these cells the sharp assumption fails on 23\% of instances, and the
stated first fall degree is wrong under Semaev's own definition.

\section{The ceiling on learned SAT guidance}\label{sec:bd-sat}
The SAT engines and encodings are described in \S\ref{sec:bo-sat}; this section records the ceiling they place on learned guidance.


The SAT survey dismissed ML-guided SAT without measuring it. A preregistered
study of 27 October measured the most any learned component could gain
inside a SAT point-decomposition oracle before anyone builds one
(\art{research/pdp_ml_sat_ceiling_20261027/README.md}; protocol frozen in
\texttt{8a85584fd}; symmetrised $\Sem4$, $m=3$, native CDCL+XOR, the
Trimoska-corpus curves, uniform targets; cost in conflicts). It tested
triage on target features or early solver telemetry, a fixed conflict
cutoff, and learned static branching. All 4{,}200 labels agree with
exhaustive sorted-triple search.
\begin{itemize}
\item \emph{The search is an enumerator.} A refutation costs $0.963$
  ($n=15$, $\ell=5$; 3{,}000 targets) and $0.998$ ($n=17$, $\ell=6$; 1{,}200
  targets) conflicts per sorted factor-base triple, every refutation within
  4.1\% of its family median (5{,}781 and 45{,}649 conflicts).
\item \emph{The prize is real.} A perfect unsatisfiability predictor,
  charged its telemetry conflicts, would raise relations per conflict by
  $5.2$--$9.0\times$.
\item \emph{The signal is not.} Test AUC is within noise of 0.5 for every
  feature group on both families and both splits; the permuted-label control
  reads 0.45--0.54 and the positive control (labels flipped with probability
  0.25) 0.72--0.79, so the learner works. The one interval that excludes
  0.5 ($n=15$ telemetry at 512 conflicts, $0.556$ $[0.518,0.599]$) includes a
  ``decided within 512 conflicts'' flag and gains $0.99\times$. Train-fitted
  gains up to $4.7\times$ fall to about $1\times$ held out.
\end{itemize}
The study closed by its preregistered rule. A related instrument finding is recorded in the atlas: five ``unverified
SAT answers'' of an earlier binary solver comparison were quadratic-twist
decompositions (\texttt{EV-ICPERF-390707};
\art{crypto-autoresearcher:knowledge/leads/ecdlp/ic-prime/LEAD-ICP-icperf-binary-pdp-solver.yaml}).

\section{F6-IC with a pair index is the meet-in-the-middle scan}\label{sec:bd-f6}
The engine itself, its gate and its other uses are described in \S\ref{sec:bo-f6}; this section records what it established for binary decomposition.


The F6-IC gate of \S\ref{sec:groebner} decided a node exactly only once one
summand's coordinate bits were all assigned. On 5 November it was changed
to run the same exact closure as soon as a summand's assigned bits leave at
most 64 base codes, and the change is now on by default
(\art{research/f6_few_candidates_20261105/README.md}; nine Koblitz cells,
5 rounds $\times$ 3 repetitions, L0).
\begin{itemize}
\item Verdicts and indices are identical on all 80 target--engine pairs;
  gate pair lookups and group additions are identical on all 40 F6-IC
  targets.
\item F6-IC CPU falls $2.67$--$10.29\times$ (worst pairing of rounds
  $\ge1.66\times$), word XORs $1.84$--$10.13\times$, and 204--1{,}049
  reductions per cell become 5--16; the inherited-F4 control reads
  $0.87$--$1.05\times$.
\end{itemize}
Because the pair lookups and group additions do not change, the algebra was
not ruling out any candidate $x$-coordinate on these cells: with the pair
index, F6-IC works as the pair-index meet-in-the-middle scan, and the
closure removes the algebra around it. The source marks the inference that
it cannot beat ecbench's \texttt{mitm} oracle, being that oracle in the
limit, as unverified. Two other items in the same study changed nothing:
Four Russians batching of displaced rows would spend about 32 row
operations per pivot on tables against about 1{,}500 row XORs per node it
replaces, and a parallel split of one target's tree gives $1.59$--$2.88\times$
wall on cells with no decomposition and $0.94$--$1.11\times$ where most
targets decompose (opt-in, refused by ecbench).

The same conclusion holds from the cost side. A desk evaluation of 28
candidate point-decomposition strategies against the lead atlas found 17
that collide with existing leads, 7 partial overlaps and 4 with no match,
rated least promising
(\art{research/pdp_strategy_evaluation_20261031/README.md}). On the ladder
rows the walk is 61--82\% of $S$ and the matrix at most 0.10\%, and on the
committed algebraic-oracle rows the oracle's share of $S$ is exactly zero
when native operations are unpriced, so an algebraic-oracle lever has a
ceiling of $1.00\times$ there.

\chapter{ECC2K-130, ECC2-131, factor-base selection, composite degree and GHS}\label{ch:binarytargets}

The named binary targets are in \S\ref{sec:ecc2k130} (ECC2K-130) and
\S\ref{sec:binary} (ECC2-131), with the priced routes of
Table~\ref{tab:prog_koblitz-ecc2k130}. This chapter adds what those sections
do not have:
\begin{itemize}
\item two independent budget tables and the isogeny-class census from the
  \texttt{cryptanalysis} repository, the Torelli toy ladder, and the state of
  the modular carrier search (\S\ref{sec:bt-ecc2k130});
\item the ECC2-131 symmetry-breaking and stage measurements
  (\S\ref{sec:bt-ecc2-131});
\item factor-base selection by a structural law and by a learned cost
  model, with the exact standing of the one validated claim
  (\S\ref{sec:bt-fbselect});
\item quasi-subfield polynomials at $n=131$ (\S\ref{sec:bt-qsp});
\item composite-degree curves (\S\ref{sec:bt-composite});
\item binary GHS and Weil-restriction screens outside the late-October cubic
  catalogue (\S\ref{sec:bt-ghs}).
\end{itemize}
The \texttt{cryptanalysis} sources are read at the head of the branch
\texttt{claude/cryptanalysis-repo-setup-p59kek} (22 October), not a
\texttt{main} branch. None of the studies here carries a validation file in
a repository read for this treatise.

\section{ECC2K-130 beyond the priced routes}\label{sec:bt-ecc2k130}

\paragraph{What one decomposition may cost.} Two \texttt{cryptanalysis}
studies turned the ECC2K-130 cost model around: fix rho at $2^{60.81}$
iterations on the $\langle-1,\tau\rangle$ classes, count one group addition
as one rho iteration, and ask how much one decomposition may cost if the
whole attack is to tie rho. The subspace version
(\art{cryptanalysis:experiments/pdp-scaling/README.md}) finds no budget at
$m=3$ (linear algebra alone exceeds rho), $2^{12.2}$ iterations at $m=4$,
$\ell=29$ (about 0.5\,ms), $2^{32.8}$ at $m=5$, $\ell=28$ (about 12 minutes)
and $2^{36.8}$ at $m=6$, $\ell=24$. Against these, its meet-in-the-middle
decomposition costs $2^{58}$, $2^{84}$ and $2^{72}$ group additions per
point, gaps of $+46$, $+51$ and $+35$ bits, and the fitted algebraic costs
(SAT, F4) are near $2^{180}$ at $m=4$; SAT grows about 4 ($m=3$) and 7 ($m=4$)
bits per extra dimension of the base. The Frobenius-quotient version
(\art{cryptanalysis:experiments/frobenius-quotient-m4/README.md}) uses the
one quotient that remains at $n=131$, the orbit quotient of the non-linear
weight base $F_w=\{P: 1\le\mathrm{HW}(x(P))\le w\}$ in a normal basis
(\S\ref{sec:kc-nonlinear}), with columns $B/131$ credited in full.
Table~\ref{tab:bt-fq} gives its best weight per arity.

\begin{table}[t]
\centering\small
\begin{tabular}{@{}rrrrrrr@{}}
\toprule
$m$ & $w$ & $\log_2B$ & $\log_2$ budget/attempt & required exponent in $B$ & MITM exponent & MITM gap, bits \\
\midrule
2 & --- & & \multicolumn{4}{l}{none: linear algebra alone exceeds rho at every $w$} \\
3 & 6 & 31.6 & 0.5 & 0.01 & 2 & $+65.0$ \\
4 & 6 & 31.6 & 31.1 & 0.98 & 2 & $+34.4$ \\
5 & 5 & 27.2 & 40.6 & 1.49 & 3 & $+44.0$ \\
6 & 4 & 22.5 & 45.3 & 2.01 & 3 & $+25.3$ \\
8 & 3 & 17.5 & 50.3 & 2.87 & 4 & $+23.7$ \\
12 & 2 & 12.1 & 52.9 & 4.38 & 6 & $+25.6$ \\
\bottomrule
\end{tabular}
\caption[]{ECC2K-130 on the Frobenius-orbit quotient of the weight base:
the most one decomposition may cost if the attack is to tie rho
($2^{60.81}$), and the meet-in-the-middle floor against it. Arithmetic
only (rows $m=7,9,10,11$ omitted). Source:
\art{cryptanalysis:experiments/frobenius-quotient-m4/README.md}.}
\label{tab:bt-fq}
\end{table}

The quotient raises the budget from $2^{12.2}$ to $2^{31.1}$ at $m=4$, which
is the factor 131 on relations and $131^2$ on linear algebra that rho
already collects through $\sqrt{2\cdot131}$. It does not remove the
obstruction. A decomposition that uses only the group law needs
$\abs{L_1}\abs{L_2}\ge\#E\approx2^{131}$, at least $2^{65.5}$ per attempt;
a method that uses the $x$-coordinate algebra must reach the required
exponent column, which is below $m/2$ for every $m$, against $m$
(exhaustive, WDSat, FES), $m-1$ (hybrid) and $\lceil m/2\rceil$ (MITM, floored
at $2^{65.5}$). At $m=4$ the required exponent is $0.98$: one decomposition
would have to cost less than listing the base once. The same study measured
$m=4$ on the weight base at toy size with CryptoMiniSat (chained $\Sem3$,
weight as a cardinality bound): every completed instance at $n=9$--15
returned a lifted, re-added decomposition with no spurious model and no
oracle disagreement, and the descriptive slope of median CPU on $\log_2B$ over
four uncensored cells is $1.34$.

\paragraph{The isogeny class, by level.} The class of
$E_0: y^2+xy=x^3+1$ over $\FF_{2^{131}}$ has four endomorphism levels,
conductors 1, 263, $p$ and $263p$ with $p=146{,}505{,}763{,}881{,}528{,}721$
(\art{cryptanalysis:experiments/ecc2k130-isogeny-class/README.md}):
\begin{itemize}
\item level 1 ($E_0$): native and effective rho $2^{60.81}$;
\item level 263, 262 floor curves: native rho $2^{64.33}$, effective
  $2^{60.81}$ through a verified 263-isogeny to $E_0$ costing at most
  $2^{15.6}$ $\FF_q$-multiplications (the 262 $j$-invariants independently
  rederived by V\'elu from $E_0[263]$, 262 of 262 matching);
\item level $p$, $p+1\approx2^{57.0}$ curves, and level $263p$,
  $262(p+1)\approx2^{65.1}$ curves: native $2^{64.33}$, effective in
  $[2^{60.81},2^{64.33}]$.
\end{itemize}
No curve in the class is easier than $E_0$. The $p$-level interval would
close only with a characteristic-2 higher-dimensional (Kani) isogeny
pipeline, and constructing any such curve costs about $2^{70}$
$\FF_q$-multiplications, more than solving natively on it. The treatise's
262-curve figure and the $11.45\times$ negation-only penalty of
\S\ref{sec:ecc2k130} are the level-263 row.

\paragraph{A Jacobian carrier: the toy ladder.} \S\ref{sec:ecc2k130} leaves
one hypothetical route open, a genus-130 Jacobian isogenous to the simple
130-dimensional abelian variety that carries $\langle G\rangle$, and quotes
the Torelli codimension 8{,}128 against it. A frozen toy test of that rarity
heuristic (\art{research/ic_next10_20261014/10-torelli-toy/README.md})
enumerated curves over $\FF_2$ whose Frobenius polynomial is divisible by
the trace-zero part $A_n$ of $\mathrm{Res}_{\FF_{2^n}/\FF_2}(E_a)$:
\begin{itemize}
\item \emph{$n=5$: negative, with low power.} No genus-5 Jacobian
  (hyperelliptic and trigonal completely, canonical curves by a uniform
  sample of $2^{21}$ nets), and no plane quintic or hyperelliptic curve of
  genus 6--8, carries $A_5$; but the heuristic predicts only about 0.42
  classes per sign at genus 5.
\item \emph{$n=7$: the literal reopen criterion is met, by structure $n=131$
  lacks.} 1{,}536 hyperelliptic models of genus 7 (mass 0.5) per sign have
  Jacobians isogenous to the full Weil restriction, against a prediction of
  $2^{-14}$--$2^{-18}$. Here $\operatorname{ord}_7(2)=3$, $A_7$ splits into
  two 3-dimensional factors, and genus $7=2^3-1$ is what GHS gives with
  magic number 4 (the attribution to GHS is the author's interpretation).
\item \emph{Calibration.} The heuristic is right to within
  $2^{-3.5}$--$2^{+1.2}$ for the $A_3$ control at excess genus at most 3,
  and under-predicts every structured class that occurs, by up to $2^{19}$
  at $n=7$.
\end{itemize}
So the ECC2K-130 rarity argument rests on the absence of such structure at
$n=131$ (no small magic numbers, a simple $A$), not on the number 8{,}128.
The atlas lists $n=11$ and 13 as the next rungs
(\art{crypto-autoresearcher:knowledge/leads/ecdlp/ic-binary/LEAD-ICB-ecc2k130-jacobian-carrier.yaml}).

\paragraph{The modular carrier search.} A separate line looks for a
low-genus or modular carrier of the relevant CM factor through an explicit
modular bridge (\art{research/koblitz_modular_ic_20261018/README.md}). It
has run 26 rounds of exact constructions with Lean identities and no
index-calculus rank. Round 22 computed the integral equivariant recovery on
norm-zero targets (eight Smith factors 1, kernel rank 24); round 26 shows by
exact diamond norms that the invariant quotient does not carry the intended
CM factor under the retained bridge, and records the next Riemann--Roch
space (degree 128, dimension 68 under the genus-61 bridge) whose basis is
still required.

\paragraph{The ranking of 23 October.} A leads survey of that date ranked
six degree and closure leads and three rho-accounting items by planning
cost against best-case effect, both judgements and not measurements
(\art{research/ecc2k130_leads_survey_20261023/report.typ}). Its first lead,
the matched null at $(22,7)$, ran on 28 October (\S\ref{sec:bd-guess}). It
quotes the unconverted distances $+7.62$ bits ($m=4$, degree 4), $+10.32$
(degree 5) and a free-oracle floor of $2^{56.46}$, and no lead moves a priced
row below rho.

\section{ECC2-131}\label{sec:bt-ecc2-131}

ECC2-131 is a random binary curve over $\FF_{2^{131}}$ ($z^{131}+z^{13}+z^2+z+1$,
cofactor 2, $a\notin\FF_2$, so no Frobenius saving), and 131 is prime, so
there is no Weil descent to a subfield. \S\ref{sec:binary} gives its program,
gates and degree falsifiers. Two measurements were added.

\paragraph{Symmetry breaking and double large primes at $m=2$.} With the
base $V_\ell=\{x:\deg x<\ell\}$ and a large-prime domain $V_{\ell'}$, the
binary $\Sem3$ linearises in $e_1=x_1+x_2$ to $t^2+t=e_1^2+b/x_R^2$, and
because $V_{\ell'}$ is closed under addition the symmetric oracle can sweep
$e_1$ (each unordered pair once, one batched inversion) instead of lifting
each $x_1$ with its own inversion
(\art{research/certicom131_symlp_20261025/REPORT.typ}; single thread, seed
20261025; IC step and rho step in the same field code and process).
Table~\ref{tab:bt-symlp} gives the step rates. For ECCp-131 the interval
$[0,B)$ is not closed under addition and each step needs two square roots,
so the symmetric oracle is slower than its baseline. Direct relations (both
abscissae in $V_\ell$) are about $5\cdot10^{-7}$ of decompositions on the
binary calibration curves; the two-large-prime graph is what makes the bulk
usable, and the planted pipeline recovers a logarithm through it on both
real 131-bit groups.

\begin{table}[t]
\centering\small
\begin{tabular}{@{}llrrrrr@{}}
\toprule
Family & Curve & $n$ & base M/s & sym M/s & sym/base & IC/rho step rate \\
\midrule
binary & ECC2-131 & 131 & 0.18 & 5.24 & $28.7\times$ & 1.20 \\
binary & calibration & 127 & 0.36 & 7.59 & $20.8\times$ & 1.74 \\
binary & calibration & 61 & 0.39 & 4.52 & $11.6\times$ & 1.04 \\
prime & ECCp-131 & 131 & 0.17 & 0.09 & $0.6\times$ & 0.02 \\
prime & calibration & 61 & 0.37 & 0.19 & $0.5\times$ & 0.03 \\
\bottomrule
\end{tabular}
\caption[]{Decomposition-oracle step rates at $m=2$, single thread: the
per-point baseline, the symmetric $e_1$ oracle, and the symmetric step rate
over the rho adding-walk step rate on the same host. A step-rate ratio, not
a cost ratio. Source:
\art{research/certicom131_symlp_20261025/REPORT.typ}.}
\label{tab:bt-symlp}
\end{table}

The study's derived cost model (every row computed, not measured) puts an
$m=2$ solve at about $2^{2n/3}$ decomposition steps at the balancing
dimension $d=(n+1)/3$, which agrees with the merged algebraic ladder's
measured trial slope $0.647$ $[0.639,0.655]$, against rho's
$\sqrt{\pi r/4}$: $\log_2$ steps $87.3$ against $64.8$ for ECC2-131 and $65.3$
for ECCp-131, a distance of about $2^{22}$ steps before the linear-algebra
floor, which the program gates place above rho for $m\le3$. Both
accelerators are constants on the $2n/3$ exponent; they would matter at
$m\ge4$, which the study does not reach.

\paragraph{Stage rates.} On one isolated sibling pair of an Intel Xeon
6975P-C host, five samples per cell
(\art{research/certicom_131_stage_bench_20261023/REPORT.md}):
\begin{itemize}
\item rho state steps per second: 207{,}822 (ECC2K-130), 207{,}602
  (ECC2-131), 87{,}055 (ECCp-131);
\item full-field $\Sem3$ quadratic-root queries per second: 48{,}080,
  48{,}186 and 9{,}522.
\end{itemize}
The two binary curves share the field polynomial and differ by $0.11\%$ in
point updates; ECCp-131's generic rho step is $0.419\times$ and its root
query $0.198\times$ ECC2K-130's. These are generic-binary CPU stage ratios,
not solve costs; the optimised clients of \S\ref{sec:ecc2k130} and
\S\ref{sec:binary} walk 8.5--13.2 million steps per second.

\section{Factor-base selection}\label{sec:bt-fbselect}

Named subspace properties (greedy coverage, cosets, geometric progressions,
normal against polynomial bases) are exhausted at scale
(\S\ref{sec:binary}, \S\ref{sec:frontier}). Two further results choose the
base by the solving degree rather than by the yield.

\paragraph{A structural law for $m=2$.} Over 656 factor bases (7 families,
9 fields) on $y^2+xy=x^3+1$, the closure solving degree of the two-summand
system is
\[
D=2+[e<-3.5]+[e<-8.5],\qquad e=n-\dim V-\dim V^{(2)},
\]
with $V^{(2)}$ the span of pairwise products, at mean absolute error 0.06
degrees, 0.08 with each field size held out
(\art{cryptanalysis:experiments/pdp-degree-heuristics/README.md}). For prime
$n$, $\dim V^{(2)}\ge2\ell-1$, and below $n-1$ only geometric progressions
$c\cdot\mathrm{span}\{1,g,\dots,g^{\ell-1}\}$ attain it; the literature's
$\mathrm{span}\{1,z,\dots,z^{\ell-1}\}$ is one. Where the rule puts a random
base a degree higher, the progressions are 15--40 times cheaper per query
and per relation on average, up to 100; elsewhere 1.5--2 times; the yield is
the same. Two further findings matter for pricing: the homogeneous degree of
regularity of these systems is exactly $\ell+1$, because the two-summand
top parts are bilinear across the variable blocks, while the solving degree
is 2--4, so the semi-regular prediction misses by 0.5--1.2 degrees on
average; and at $m=3$ the formulation decides whether the base matters (in
the symmetric formulation the progression is 15--40$\times$ cheaper per query
at $\ell=3$ at equal degree, and 500--1{,}500$\times$ where the random base
needs one more degree).

\paragraph{A learned cost model.} The lead
(\art{crypto-autoresearcher:knowledge/leads/ecdlp/ic-binary/LEAD-ICB-learned-cost-model-factor-base-selection.yaml},
last reviewed 6 November) fits a gated log-linear (ridge) model of
per-relation cost on structural features of $V$ (predicted exact yield,
solving degree, trace and product-space statistics) and compares its pick
from a pool with the hand rule ``lowest expected solving degree first''. The
label producer is the structural-law experiment above. As the atlas records
it:
\begin{itemize}
\item \emph{Stage.} $1.82$ bits per relation $[1.36,2.26]$ better than the
  hand rule at $\ell\le4$ (13 cells), $1.13$ $[0.45,1.62]$ at $\ell=5$.
\item \emph{Across curves, frozen, no refit.} $+1.84$ $[0.95,2.68]$ on
  $a_2=1$ Koblitz curves; $+0.92$ $[0.54,1.34]$ on ordinary curves, not
  confirmed on 16 fresh ordinary curves at $\ell=4$ ($+0.43$ $[0.15,0.74]$
  against a $1.08$ threshold), where the hand rule already sits near the pool
  best.
\item \emph{End to end.} $0.93$ bits $[0.36,1.48]$ less counted cost per
  verified logarithm over 9 cells ($n\le41$). At $\ell=4$ on ordinary
  curves, $n=19$, the complete-logarithm check is a null with the sign
  reversed: $-0.45$ $[-0.98,+0.04]$ bits at $w=0$.
\item \emph{Native ecbench on the isolated EKS route}, 5 November: K1
  ($a_2=1$) Koblitz curves, $\ell=3$, \texttt{ic.pipeline} with explicit
  binary-subspace bases, 4 targets $\times$ 5 rounds per size, 192 of 192
  records verified. $\log_2(S_{\text{rule}}/S_{\text{ridge}})$ in counted
  complete-logarithm cost is $+1.83/+0.80/+0.80$ bits at $n=17$ and
  $+0.96/+2.47/+2.47$ at $n=19$ for $w=0/0.1/1$, every interval's lower end
  above 0. The ridge pick stays behind \texttt{rho.signed\_frobenius\_strong}
  by $1.8/14.8/17.7$ bits ($n=17$) and $3.3/15.1/17.9$ ($n=19$). The
  producer's first ``about 9 bits, identical across the band'' was a
  pinned-calibration point and was corrected.
\end{itemize}
\emph{Standing, exactly.} The atlas lead on crypto-autoresearcher
\texttt{main} (commit \texttt{8d61579904}, 6 November) records the EKS
result as ``independently validated 2026-11-06 with VALIDATION.json, 160 of
160 replays reproduced on a second CPU''. The experiment, its protocol, its
validator report and that \texttt{VALIDATION.json} live in the external
repository \texttt{aburan28/ml-cryptanalysis}
(\texttt{experiments/fb-cost-model-curves-eks}), which this session could
not open; the same atlas entry's evidence note says the validation is ``in
a follow-up PR''. Nothing about the study is on \texttt{main} of
\texttt{crypto} or \texttt{crypto-autoresearcher} except the atlas text. So
for this treatise the validated claim is an unverified assertion recorded
by the atlas: if it holds, it is the family's only validated, measured
improvement over a binary index-calculus baseline, about 1--2.5 bits at toy
size, $n=17$, 19, with its scaling in $n$ not yet measured (K1 $n=23$ was
frozen on 6 November). It moves no row toward rho by more than the
1.8--17.9 bits it leaves.

\section{Quasi-subfield polynomials at $n=131$}\label{sec:bt-qsp}

\S\ref{sec:binary} records the $\FF_{2^7}$ census that showed the HKPYY
Lemma~4.1 bound~\cite{hkyy2018} tight. The cost side is closed by theorem
(\art{crypto-autoresearcher:analysis/qsp-ecc2k130/README.md}; the atlas
records it as a theorem from Euler and Petit's work):
\begin{itemize}
\item every linearised quasi-subfield polynomial has $\beta\ge3/4$, so under
  the HKPYY cost model the exponent is at least $0.93$, at least $2^{122}$ at
  $n=131$ against rho's $2^{60.9}$;
\item with $\FF_2$ coefficients only $X^2+X$ ($n'=1$) and $L_{\Phi_{131}}$
  ($n'=130$, $\beta=0.99994$) exist at $n=131$;
\item multiplicative quasi-subfield polynomials give $\abs V=264$ or $p+1$,
  with $\beta\ge1.02$;
\item the cost model needs $\beta\approx0.0958$ to reach $2^{60.9}$.
\end{itemize}
At toy size the yield is no better than random: at $n=31$, $k=15$, the QSP
base $t^{15}+t^7+t^3+t+1$ ($\beta=0.964$) yields $0.214$ $(0.190$--$0.240)$
against $0.204$--$0.230$ for random subspaces, and the route's kill criterion
fired (\art{research/ic_next10_20261014/08-quasi-subfield/README.md}).

\section{Composite degree}\label{sec:bt-composite}

Composite $n$ is the one binary setting where index calculus can be priced
below rho on time alone. Chapter~\ref{ch:obstructions} records the lift-dense
twisted-line base on c2pnb208w1 and c2pnb368w1 ($+0.50$ bits, modelled), and
\S\ref{sec:transfer} the GHS reachability of the Oakley curves.

\paragraph{c2pnb: Gaudry--Diem $n=4$, as a cost model.} The five ANSI X9.62
example curves c2pnb176v1, 208w1, 272w1, 304w1 and 368w1 have coefficients in
$\FF_{2^{16}}$ and live over $\FF_{Q^4}$, $Q=2^{4k}$, $k=11,13,17,19,23$. A
zero-experiment derivation
(\art{crypto-autoresearcher:analysis/ecdlp2m-iteration-20261009/c1-gaudry-n4-c2pnb/README.md})
prices Gaudry--Diem $n=4$~\cite{gaudry2009,diem2011} with a base of points
with $x\in\FF_Q$ reduced modulo negation and $\pi_{16}$ ($N=Q/(2k)$ orbits)
and double large primes~\cite{gtthd2007}, against matched rho of $78.10$,
$93.98$, $125.78$, $141.71$ and $173.57$ bits. Under three metrics:
\begin{itemize}
\item \emph{time only:} below rho on all five rows for every per-decomposition
  cost $c_{\text{trial}}\in[2^{15.5},2^{32.5}]$, by $+1.46$, $+5.80$, $+14.19$,
  $+18.34$ and $+26.60$ bits at $c_{\text{trial}}=2^{30}$;
\item \emph{time $\times$ memory against van Oorschot--Wiener:} worse on every
  row, by 35--56 bits at $c_{\text{trial}}=2^{30}$;
\item \emph{mesh area-time:} better only at 304w1 ($+1.28$) and 368w1
  ($+0.97$), and only for $c_{\text{trial}}\le2^{27}$.
\end{itemize}
The verdict turns on $c_{\text{trial}}$, never measured at toy size, the
per-worker FGLM memory ($D^2$ with $D=4096$, brief-supplied) and the
linear-algebra constant; a toy ladder over $\FF_{2^4}$ lifted to $\FF_{2^{4k}}$,
$k\in\{5,7,9,11,13\}$, about $2\cdot10^4$ decompositions, would pin all three
and has not run. The atlas adds that the $\FF_{2^{16}}$-rational base is
structurally empty in $G$ and that a naive trace-zero extrapolation is a
regime-invalid trap
(\art{crypto-autoresearcher:knowledge/leads/ecdlp/ic-binary/LEAD-ICB-composite-degree-c2pnb.yaml}).

\paragraph{Oakley groups 3 and 4: $n=5$.} The twisted line exists on both
Oakley EC2N curves: a $u$ with $T(u)=T(c/u)=0$ was verified in $\FF_{2^N}$
($N=155$ and 185, $n=5$), and all 200 sampled elements of $u\FF_q$ lift
against 104 and 99 of 200 for the standard line, so as a Gaudry--Diem $n=5$
base each decomposition trial would have $2^5=32$ times the relations of the
standard base; on c2pnb176v1, 272w1 and 304w1 ($\operatorname{Tr}(a)=1$)
none of 200 lift (\art{research/oakley_lift_dense_20261027/REPORT.typ}). This
is a per-trial yield statement. Three approved Oakley contracts have no
runs; the atlas prices the GHS genus-16 transfer at $+1.1$ bits (group 3) and
$+10.9$ bits (group 4) of advantage after the 11 October magic-number
correction, by formula
(\art{crypto-autoresearcher:knowledge/leads/ecdlp/ic-binary/LEAD-ICB-oakley-quintic-tower.yaml}).

\section{Binary GHS and Weil-restriction screens}\label{sec:bt-ghs}

The late-October cubic GHS catalogue over $\FF_4$, $\FF_8$ and $\FF_{16}$
belongs to the transfer record (\S\ref{sec:transfer}). Four binary screens
outside it bear on index calculus.

\paragraph{Decomposition coverage at $n=2$ is $3/8$, not Poisson.} For
generic $E/\FF_{q^n}$ and the base $x\in\FF_q$, the share of points that are
a sum of $n$ base points was computed exactly over the whole group
(\art{research/weil_descent_angles_20261018/REPORT.md}). In odd
characteristic at $n=2$--5 it matches a random $x$-set of the same size
(paired residual $\abs z\le0.94$). Binary $n=2$ is the exception: coverage
tends to $3/8$, not the Poisson $1-e^{-1/2}=0.3935$ ($z=-4.72$), and the cause
is proved. Every $x\in\FF_q$ has $\operatorname{Tr}_{\FF_{q^2}/\FF_2}(x)=0$,
so the base lies in one coset of $2E$; the descended relation map is
$\FF_2$-affine with image in the hyperplane $\operatorname{Tr}(w/x_3^2)=0$; so
a point decomposes only if $\operatorname{Tr}(x_3)=\operatorname{Tr}(a)$, half
the group, and inside that half the multiplicities are exactly
$\mathrm{Bin}(2,1/2)$. Because attack targets lie in the odd-order subgroup
inside $2E$, the relevant coverage is that of $2E$, exactly $3/4$ against
$0.3935$; the study corrected its first reading (``neutral to slightly
negative'', measured over the whole group) on this point. The identity holds
on all 21 curves.

\paragraph{IPsec group 3.} An archived search over the isogeny class of RFC
2409 Oakley group 3 ($\FF_2[z]/(z^{155}+z^{62}+1)$, $N=12r$) screened 24{,}492
Frobenius orbits and found no low-genus candidate in the implemented tests;
the whole class was not enumerated and no resistance claim follows
(\art{research/ghs_ipsec_isogeny_class_20261023/README.md}).

\paragraph{Generalised binary GHS, explicitly.} Full nonzero-$c$ Hess screens
at absolute degrees 6, 8, 10 and 12 (63, 255, 1{,}023 and 4{,}095 values) gave
explicit maps and attack controls
(\art{research/ghs_generalized_20261024/REPORT.md}). At degree 12, 147
profiles have genus at most 2, including 15 of genus 1 over $\FF_{2^4}$, and the
Weil bound ($\#\mathrm{Jac}\le25$ and $\le81$) rules out a homomorphism from
the order-199 subgroup; the other 232 profiles are index-two genus-2
profiles that a subfield change maps to one canonical form, presenting $P$ as
$T-P$ for the 2-torsion point $T$. Every exact instance (subgroups of order
19, 67, 43 and 199) verified end to end, and all belong to known
GHS~\cite{ghs2002} and Hess families.

\paragraph{A degree-73 descent at $n=37$.} A separable degree-73 isogeny of
$y^2+xy=x^3+1$ over $\FF_{2^{37}}$ (74 rational cyclic kernels; a squarefree
degree-36 kernel polynomial dividing the degree-2{,}664 division polynomial)
was tested as a presentation change for the Boolean stage
(\art{research/koblitz_isogeny_descent_37_results_20260925/RESULTS.md}). All
128 systems completed and every planted decomposition verified on both
curves, but the preregistered criterion (at least a 10\% F5 matrix XOR
reduction in every cell, with no completion-degree increase) failed: the
neighbour cut XORs by at least 10\% in 4 of 8 cells, regressed up to
$2.00\times$, and raised the completion degree in 7 of 64 paired systems.
The committed raw archive was found not to be gzip and was restored from the
original workflow artifact; the result is unchanged.
